
\section{标记、线型、（背景）颜色与颜色映射}
\label{sec:markers}

以下 \Tikz{} 选项可用于绘图。


\subsection{标记}

该列表摘录自~\cite[第~29 节]{tikz}：
%
\pgfmanualpdflabel{/tikz/mark}{}%
\begingroup
\newenvironment{longdescription}[0]{%
    \begin{list}{}{%
        \leftmargin=4.7cm
        \setlength{\labelwidth}{4.7cm}%
        \renewcommand{\makelabel}[1]{\hfill\textbf{\texttt{##1}}}%
    }%
}{%
    \end{list}%
}%
\def\showit#1{%
    \tikz\draw [
        gray,
        thin,
        mark options={fill=yellow!80!black,draw=black,scale=2},
        x=0.8cm,y=0.3cm,
        #1,
    ] plot coordinates {(0,0) (1,1) (2,0) (3,1)};%
}%
\def\showitpgfplots#1{%
\begin{tikzpicture}[baseline]
    \begin{axis}[
        anchor=north,
        xticklabels=,
        yticklabels=,
        zticklabels=,
        width=5cm,
    ]
        \addplot3 [
            gray,
            thin,
            mark options={
                scale=2,
                fill=yellow!80!black,
                draw=black,
            },
            #1,
        ] coordinates {(0,0,0) (0.3,0.6,0.3) (2,0,0.1) (2.3,1,0.2)};
    \end{axis}
\end{tikzpicture}%
}%
\begin{longdescription}
    \item[mark=*] \showit{mark=*}
    \item[mark=x] \showit{mark=x}
    \item[mark=+] \showit{mark=+}
%    \item[mark=ball] \showit{mark=ball}
\end{longdescription}
%
配合 |\usetikzlibrary{plotmarks}| 后还可用：
%
\begin{longdescription}
    \item[mark=$-$] \showit{mark=-}
    \item[mark=$\vert$] \showit{mark=|}
    \item[mark=o] \showit{mark=o}
    \item[mark=asterisk] \showit{mark=asterisk}
    \item[mark=star] \showit{mark=star}
    \item[mark=10-pointed star] \showit{mark=10-pointed star}
    \item[mark=oplus] \showit{mark=oplus}
    \item[mark=oplus*] \showit{mark=oplus*}
    \item[mark=otimes] \showit{mark=otimes}
    \item[mark=otimes*] \showit{mark=otimes*}
    \item[mark=square] \showit{mark=square}
    \item[mark=square*] \showit{mark=square*}
    \item[mark=triangle] \showit{mark=triangle}
    \item[mark=triangle*] \showit{mark=triangle*}
    \item[mark=diamond] \showit{mark=diamond}
    \item[mark=diamond*] \showit{mark=diamond*}
    \item[mark=halfdiamond*] \showit{mark=halfdiamond*}
    \item[mark=halfsquare*] \showit{mark=halfsquare*}
    \item[mark=halfsquare right*] \showit{mark=halfsquare right*}
    \item[mark=halfsquare left*] \showit{mark=halfsquare left*}
    \item[mark=Mercedes star] \showit{mark=Mercedes star}
    \item[mark=Mercedes star flipped] \showit{mark=Mercedes star flipped}
    \item[mark=halfcircle] \showit{mark=halfcircle}

        其中一半以白色填充（更准确地说，是以 |mark color| 填充）。

    \item[mark=halfcircle*] \showit{mark=halfcircle*}

        其中一半以白色填充（更准确地说，是以 |mark color| 填充），
        另一半则以实际的 |fill| 颜色填充。

    \item[mark=pentagon] \showit{mark=pentagon}
    \item[mark=pentagon*] \showit{mark=pentagon*}
    \item[mark=ball] \showit{mark=ball,mark options={ball color=yellow!80!black,scale=2}}

        该标记较为特殊；若使用数量很大，很容易生成过大的输出文件。
        它的另一特殊之处在于需要设置 |ball color|（在本例中为
        |ball color=yellow!80!black|）。

    \item[mark=text] \showit{mark=text,every mark/.append style={scale=0.5}}

        该标记较为特殊，因为其样式可自由配置。所使用的字符
        （甚至整段文字）由一组变量配置，详见下文。

    \item[mark=cube] \showitpgfplots{mark=cube}

        该标记仅在 \PGFPlots{} 坐标轴内部可用；它绘制一个面与坐标轴
        平行的立方体。其尺寸可分别配置，详见下文。

    \item[mark=cube*] \showitpgfplots{mark=cube*}
    \item[用户自定义] 可以使用 |\pgfdeclareplotmark| 定义新的标记，
        详见下文。
\end{longdescription}
%
以下额外选项用于绘制上述所有标记：
%
\begin{codeexample}[code only]
\draw [
    gray,
    thin,
    mark options={
        scale=2,fill=yellow!80!black,draw=black,
    },
]
\end{codeexample}
%
请参阅第~\ref{sec:colors} 节了解如何修改 |draw| 与 |fill| 颜色。
注意，所提供的每个标记都可以通过 |mark options={rotate=90}| 或
|every mark/.append style={rotate=90}| 自由旋转。

\begin{key}{/tikz/mark size=\marg{dimension} (initially 2pt)}
    该 \Tikz{} 选项用于将标记尺寸设置为 \meta{dimension}。对于圆形标记，
    \meta{dimension} 为半径；对于其他绘图标记，则约为其宽度和高度的一半。
\end{key}

\begin{pgfplotsxykey}{cube/size \x=\marg{dimension} (initially |\textbackslash pgfplotmarksize|=2pt)}
    分别为每个坐标轴设置 |mark=cube| 的尺寸。
\end{pgfplotsxykey}

\begin{key}{/tikz/every mark}
    可以重新配置该 \Tikz{} 样式，以设置标记的外观选项，
    例如颜色或缩放、旋转等变换。\PGFPlots{} 会将其 |cycle list|
    选项附加到该样式。
    %
\begin{codeexample}[]
\begin{tikzpicture}
\begin{axis}[y=2cm]
    \addplot coordinates
        {(-2,0) (-1,1) (0,0) (1,1) (2,0)};
\end{axis}
\end{tikzpicture}
\end{codeexample}

\begin{codeexample}[]
\tikzset{every mark/.append style={scale=2}}
\begin{tikzpicture}
\begin{axis}[y=2cm]
    \addplot coordinates
        {(-2,0) (-1,1) (0,0) (1,1) (2,0)};
\end{axis}
\end{tikzpicture}
\end{codeexample}

\begin{codeexample}[]
\begin{tikzpicture}
\begin{axis}[y=2cm]
  \addplot+ [
    mark=halfcircle*,
    every mark/.append style={rotate=90}]
  coordinates
    {(-2,0) (-1,1) (0,0) (1,1) (2,0)};

  \addplot+ [
    mark=halfcircle*,
    every mark/.append style={rotate=180}]
  coordinates
    {(-2,-0.1) (-1,0.9) (0,-0.1) (1,0.9) (2,-0.1)};
\end{axis}
\end{tikzpicture}
\end{codeexample}

    注意，|every mark| 在某种程度上是静态的，它只会被求值一次。
    如果需要在 |scatter| 图中为每个标记分别着色，则需要使用
    |scatter/use mapped color|。
    %
\begin{codeexample}[]
\begin{tikzpicture}
\begin{axis}[y=2cm]
    \addplot+ [
        mark=ball,
        mark size=4pt,
        scatter,% enable scatter
        scatter src=rand,% the "color data"
        % configure individual appearances:
        scatter/use mapped color=
            {ball color=mapped color}]
    coordinates
        {(-2,0) (-1,1) (0,0) (1,1) (2,0)};
\end{axis}
\end{tikzpicture}
\end{codeexample}
\end{key}

\begin{stylekey}{/pgfplots/no markers}
    该键会覆盖 |\addplot| 之后由选项列表的 |cycle list| 设置的任何
    |mark| 值。

    如果将该样式作为参数提供给整个坐标轴，它会被附加到
    |every axis plot post|，从而即使对于包含标记的 |cycle list|
    也会禁用标记。
\end{stylekey}

\begin{key}{/tikz/mark repeat=\marg{integer} (initially empty)}
    允许仅绘制第 $n$ 个 |mark|，其中 $n$ 由 \meta{integer} 给出。
\end{key}

\begin{key}{/tikz/mark phase=\marg{integer $p$} (initially 1)}
    该选项用于控制绘制哪些标记。它主要与 \Tikz{} 选项
    |mark repeat=|$r$ 配合使用：它告诉 \tikzname{}，第一个要绘制的
    标记应为第 $p$ 个，随后是第 $(p + r)$ 个、第 $(p + 2r)$ 个，依此类推。
    %
\begin{codeexample}[]
\begin{tikzpicture}
    \begin{axis}[tiny]
        \addplot+ [scatter] {sin(deg(x))};
    \end{axis}
\end{tikzpicture}
\begin{tikzpicture}
    \begin{axis}[tiny]
        \addplot+ [scatter,
            mark repeat=3,mark phase=2]
                {sin(deg(x))};
    \end{axis}
\end{tikzpicture}
\end{codeexample}
    %
    此处 $p=1$ 表示第一个数据点（即 |\coordindex|$=0$ 的那个点）。
\end{key}

\begin{key}{/tikz/mark indices=\marg{index list} (initially empty)}
    允许仅绘制索引编号在参数列表中的标记。
\end{key}

\begin{key}{/pgf/mark color=\marg{color} (initially empty)}
    为 |halfcircle|、|halfcircle*|、|halfdiamond*| 与 |halfsquare*|
    标记定义\emph{额外的}填充颜色。空值使用 |white|（即初始配置）。
    值 |none| 会禁用该部分的填充。

    这些标记具有两种不同的填充颜色：一种与其他标记一样由 |fill|
    决定，另一种则是 |mark color|。
    %
\begin{codeexample}[]
\begin{tikzpicture}
\begin{axis}[y=2cm]
    \addplot [
        blue,
        mark color=blue!50!white,
        mark=halfcircle*
    ] coordinates {
        (-2,0) (-1,1) (0,0) (1,1) (2,0)
    };
    \addplot [
        red,
        mark color=red!50!white,
        mark=halfsquare*,
    ] coordinates {
        (-2,-0.1) (-1,0.9) (0,-0.1) (1,0.9) (2,-0.1)
    };
\end{axis}
\end{tikzpicture}
\end{codeexample}

    注意，该键需要 \PGF{} 2.10 或更高版本。
\end{key}

\begin{key}{/tikz/mark options=\marg{options}}
    将 |every mark| 重置为 \marg{options}。
\end{key}

\begin{key}{/pgf/text mark=\marg{text} (initially p)}
    修改 |mark=text| 所显示的文字。

    使用 |/pgf/text mark=m|：\pgfkeys{/pgf/text mark=m}\showit{mark=text,every
    mark/.append style={scale=0.5}}

    使用 |/pgf/text mark=A|：\pgfkeys{/pgf/text mark=A}\showit{mark=text,every
    mark/.append style={scale=0.5}}

    字符数量等没有任何限制。事实上，任何 \TeX{} 内容都可以作为
    \meta{text} 插入，包括图片。
\end{key}

\begin{key}{/pgf/text mark style=\marg{options for \texttt{mark=text}}}
    定义一组用于控制 |mark=text| 外观的选项。

    若 |/pgf/text mark as node=false|（默认值），则 \meta{options} 作为
    参数提供给 |\pgftext|——它仅提供 |left|、|right|、|top|、|bottom|、
    |base| 与 |rotate| 等基本键。

    若 |/pgf/text mark as node=true|，则 \meta{options} 作为参数提供给
    |\node|。这意味着你可以提供一组功能非常强大的选项，
    包括 |anchor|、|scale|、|fill|、|draw|、|rounded corners| 等。
\end{key}

\begin{key}{/pgf/text mark as node=\mchoice{true,false} (initially false)}
    配置 |mark=text| 的绘制方式：作为 |\node| 或作为 |\pgftext|。

    第一种方式非常灵活，但可能较慢；第二种方式非常快，
    通常已足够。
\end{key}

\begin{command}{\pgfdeclareplotmark\marg{plot mark name}\marg{code}}
    定义一个名为 \meta{plot mark name} 的新标记。每当使用该标记时，
    就会执行 \meta{code}。其中应包含（最好是 \PGF{} 基础层的）
    绘图命令。在 \meta{code} 执行期间，坐标系的原点表示标记应放置的
    坐标位置。

    更详细的信息请参阅~\cite{tikz} 中的 ``Mark Plot Handler'' 一节。
\end{command}

\begin{stylekey}{/pgfplots/every axis plot post (initially empty)}
    %\todosp{shall the curly brackets after ``initially'' be shown? Also
    %this style is again defined in Section~4.18.1 (see corresponding index
    %entry). Is that intended?}
    % The replication makes sense imho, but the curly braces appear to
    % be a tech debt of the manual styles
    |every axis plot post| 样式可用于覆盖为绘图分配的部分（或全部）
    绘制样式。
    %
\begin{codeexample}[]
% Overwrite any cycle list:
\pgfplotsset{
    every axis plot post/.append style={
        mark=triangle,
        every mark/.append style={rotate=90},
    },
}
\begin{tikzpicture}
    \begin{axis}[y=2cm]
        \addplot coordinates {
            (-2,0) (-1,1) (0,0) (1,1) (2,0)
        };
    \end{axis}
\end{tikzpicture}
\end{codeexample}
\end{stylekey}

标记路径与图中其他部分不同，不会被裁剪。标记要么完整绘制，要么完全不绘制。

\Tikz{} 还提供了更多用于精细调整标记的选项，详情请参阅~\cite{tikz}。


\subsection{线型}

\def\showit#1{%
    \tikz\draw[
        black,
        x=0.8cm,y=0.3cm,
        #1]
    plot coordinates {(0,0) (1,1) (2,0) (3,1)};
}%
%
以下线型在 \Tikz{} 中预定义。
%
\begin{stylekey}{/tikz/solid}
    \showit{style=solid}
\end{stylekey}

\begin{stylekey}{/tikz/dotted}
    \showit{style=dotted}
\end{stylekey}

\begin{stylekey}{/tikz/densely dotted}
    \showit{style=densely dotted}
\end{stylekey}

\begin{stylekey}{/tikz/loosely dotted}
    \showit{style=loosely dotted}
\end{stylekey}

\begin{stylekey}{/tikz/dashed}
    \showit{style=dashed}
\end{stylekey}

\begin{stylekey}{/tikz/densely dashed}
    \showit{style=densely dashed}
\end{stylekey}

\begin{stylekey}{/tikz/loosely dashed}
    \showit{style=loosely dashed}
\end{stylekey}

\begin{stylekey}{/tikz/dashdotted}
    \showit{style=dashdotted}
\end{stylekey}

\begin{stylekey}{/tikz/densely dashdotted}
    \showit{style=densely dashdotted}
\end{stylekey}

\begin{stylekey}{/tikz/loosely dashdotted}
    \showit{style=loosely dashdotted}
\end{stylekey}

\begin{stylekey}{/tikz/dashdotdotted}
    \showit{style=dashdotdotted}
\end{stylekey}

\begin{stylekey}{/tikz/densely dashdotdotted}
    \showit{style=densely dashdotdotted}
\end{stylekey}

\begin{stylekey}{/tikz/loosely dashdotdotted}
    \showit{style=loosely dashdotdotted}
\end{stylekey}
%
\noindent 由于这些样式同样适用于标记，你可能希望在标记样式中考虑使用
%
\begin{codeexample}[code only]
\pgfplotsset{
    every mark/.append style={solid},
}
\end{codeexample}
%
\noindent 。

除线型外，\PGF{} 还提供了（大量）箭头。详情请参阅~\cite{tikz}。
%
\endgroup


\subsection{边及其参数}

当 \PGFPlots{} 连接数据点时，它依赖 \PGF{} 的绘图参数来生成合适的边
（仅在 |every patch| 样式中才会修改这些参数）。

有时可能需要修改这些参数：

\begin{keylist}{
    /tikz/line cap=\mchoice{round,rect,butt} (initially butt),
    /tikz/line join=\mchoice{round,bevel,miter} (initially miter),
    /tikz/miter limit=\meta{factor} (initially 10)%
}
    这些键控制线条在拐角处如何连接。其说明超出了本手册的范围，
    感兴趣的读者请参阅~\cite{tikz}。

    下面仅举一例，说明为什么研究这些参数可能是有意义的：
% \usetikzlibrary{spy}
\begin{codeexample}[]
% requires \usetikzlibrary{spy}
\begin{tikzpicture}[spy using outlines=
    {circle, magnification=6, connect spies}]
\begin{axis}[no markers,grid=major,
    every axis plot post/.append style={thick}]
\addplot  coordinates
 {(0, 0) (0, 0.9) (1, 0.9) (2, 1) (3, 0.9) (80, 0)};
\addplot+ [line join=round] coordinates
 {(0, 0) (0, 0.9) (2, 0.9) (3, 1) (4, 0.9) (80, 0)};
\addplot+ [line join=bevel] coordinates
 {(0, 0) (0, 0.9) (3, 0.9) (4, 1) (5, 0.9) (80, 0)};
\addplot+ [miter limit=5] coordinates
 {(0, 0) (0, 0.9) (4, 0.9) (5, 1) (6, 0.9) (80, 0)};

    \coordinate (spypoint)     at (3,1);
    \coordinate (magnifyglass) at (60,0.7);
\end{axis}

\spy [blue, size=2.5cm] on (spypoint)
    in node[fill=white] at (magnifyglass);
\end{tikzpicture}
\end{codeexample}
\end{keylist}


\subsection{字号与线宽}

通常我们希望修改绘图的线宽和字号。这可以使用 \Tikz{} 的以下选项完成。

\begin{key}{/tikz/font=\marg{font name} (initially \textbackslash normalfont)}
    设置用于节点中文字的字体（例如刻度标签、
    图例或说明文字）。

    字体可以是任何 \LaTeX{} 参数，例如 |\footnotesize| 或
    |\small\bfseries|。\footnote{当然，Con\TeX{}t 和 plain \TeX{} 用户需要提供其他语句。}

    仅为特定的坐标轴描述修改字体可能会很有用，例如使用
\begin{codeexample}[code only]
\pgfplotsset{
    tick label style={font=\small},
    label style={font=\small},
    legend style={font=\footnotesize},
}
\end{codeexample}

    另请参阅第~\ref{sec:scaling:styles} 节中预定义的样式
    |normalsize|、|small| 与 |footnotesize|。
\end{key}

\begin{key}{/tikz/line width=\marg{dimension} (initially 0.4pt)}
    设置线宽。请注意，刻度线和网格线的线宽是预定义的，
    因此可能需要覆盖样式 |every tick| 与 |every axis grid|。

    |line width| 键在 \Tikz{} 中经常被修改。你应使用
\begin{codeexample}[code only]
\pgfplotsset{every axis/.append style={line width=1pt}}
\end{codeexample}
    %
    或
    %
\begin{codeexample}[code only]
\pgfplotsset{every axis/.append style={thick}}
\end{codeexample}
    %
    来修改整体线宽。若还要调整刻度线和网格线，可以使用
    %
\begin{codeexample}[code only]
\pgfplotsset{every axis/.append style={
    line width=1pt,
    tick style={line width=0.6pt}}}
\end{codeexample}
    %
    或类似样式
    %
\begin{codeexample}[code only]
\pgfplotsset{every axis/.append style={
    thick,
    tick style={semithick}}}
\end{codeexample}
    %
    `|every axis plot|' 样式可用于仅修改绘图的线宽。
\end{key}

\begin{keylist}[/tikz]{thin,ultra thin,very thin,semithick,thick,very thick,ultra thick}
    这些 \Tikz{} 样式提供不同的预定义线宽。
\end{keylist}

\begin{codeexample}[]
\pgfplotsset{every axis/.append style={
    font=\large,
    line width=1pt,
    tick style={line width=0.8pt}}}
\begin{tikzpicture}
\begin{loglogaxis}[
    legend pos=south west,
    xlabel=\textsc{Dof},
    ylabel=$L_2$ Error
]
\addplot coordinates {
    (5,8.312e-02)    (17,2.547e-02)   (49,7.407e-03)
    (129,2.102e-03)  (321,5.874e-04)  (769,1.623e-04)
    (1793,4.442e-05) (4097,1.207e-05) (9217,3.261e-06)
};

\addplot coordinates{
    (7,8.472e-02)    (31,3.044e-02)   (111,1.022e-02)
    (351,3.303e-03)  (1023,1.039e-03) (2815,3.196e-04)
    (7423,9.658e-05) (18943,2.873e-05)
    (47103,8.437e-06)};

\addplot coordinates{
    (9,7.881e-02)     (49,3.243e-02)   (209,1.232e-02)
    (769,4.454e-03)   (2561,1.551e-03)
    (7937,5.236e-04)  (23297,1.723e-04)
    (65537,5.545e-05) (178177,1.751e-05)};

\addplot coordinates{
    (11,6.887e-02)    (71,3.177e-02)   (351,1.341e-02)
    (1471,5.334e-03)  (5503,2.027e-03)
    (18943,7.415e-04) (61183,2.628e-04)
    (187903,9.063e-05) (553983,3.053e-05)};

\addplot coordinates{
    (13,5.755e-02)    (97,2.925e-02)   (545,1.351e-02)
    (2561,5.842e-03)  (10625,2.397e-03)
    (40193,9.414e-04) (141569,3.564e-04)
    (471041,1.308e-04) (1496065,4.670e-05)};
\legend{$d=2$,$d=3$,$d=4$,$d=5$,$d=6$}
\end{loglogaxis}
\end{tikzpicture}
\end{codeexample}
%
\refstepcounter{pgfmanualentry}\label{page:plotcoords:src}%
前一个示例定义的数据在本手册中多次使用；它由
\declareandlabel{\plotcoords} 引用。%

\begin{codeexample}[]
\pgfplotsset{every axis/.append style={
    font=\footnotesize,
    thin,
    tick style={ultra thin}},
}
\begin{tikzpicture}
    \begin{loglogaxis}[
        xlabel=\textsc{Dof},
        ylabel=$L_2$ Error
    ]
        % see above for this macro:
        \plotcoords
        \legend{$d=2$,$d=3$,$d=4$,$d=5$,$d=6$}
    \end{loglogaxis}
\end{tikzpicture}
\end{codeexample}


\subsection{颜色}
\label{sec:colors}

{
\def\showcolorandname#1{%
    \showcolor{#1}~\texttt{\pgfmanualpdflabel{#1}{#1}}%
}
\def\showcolor#1{%
    \tikz \draw[black,fill={#1}] (0,0) rectangle (1em,0.6em);%
}

\PGF{} 使用 |xcolor| 的颜色支持。因此，关于如何指定颜色的主要参考是
|xcolor| 手册~\cite{xcolor}。\PGF{} 手册~\cite{tikz} 则是如何为特定用途
（如绘制、填充、底纹、图案等）选择颜色的参考。本节概述~\cite{xcolor}
中颜色的指定方式（不仅限于 \PGFPlots{}）。

|xcolor| 宏包定义了一组预定义颜色，即
%
\showcolorandname{red},
\showcolorandname{green},
\showcolorandname{blue},
\showcolorandname{cyan},
\showcolorandname{magenta},
\showcolorandname{yellow},
\showcolorandname{black},
\showcolorandname{gray},
\showcolorandname{white},
\showcolorandname{darkgray},
\showcolorandname{lightgray},
\showcolorandname{brown},
\showcolorandname{lime},
\showcolorandname{olive},
\showcolorandname{orange},
\showcolorandname{pink},
\showcolorandname{purple},
\showcolorandname{teal},
\showcolorandname{violet}.

\begin{codeexample}[]
\begin{tikzpicture}
\begin{axis}[enlarge x limits=false]
    \addplot [red,samples=500] {sin(deg(x))};
    \addplot [orange,samples=7] {sin(deg(x))};
    \addplot [teal,const plot,samples=14]
        {sin(deg(x))};
\end{axis}
\end{tikzpicture}
\end{codeexample}

除预定义颜色外，还可以\emph{混合}两种（或多种）颜色。
例如，\showcolorandname{red!30!white} 包含 $30\%$ 的
\showcolorandname{red} 和 $70\%$ 的 \showcolorandname{white}。
因此，可以构造 \showcolorandname{red!70!white} 得到 $70\%$ 红色和
$30\%$ 白色，或 \showcolorandname{red!10!white} 得到 $10\%$ 红色和
$90\%$ 白色。这种混合可以用于任何颜色，例如
\showcolorandname{red!50!green}、\showcolorandname{blue!50!yellow}
或 \showcolorandname{green!60!black}。

另一种颜色混合方式也受支持，它允许取\emph{每种}分量的 $100\%$。
例如，\showcolorandname{rgb,2:red,1;green,1} 将加入 $1/2$ 份
\showcolorandname{red} 和 $1/2$ 份 \showcolorandname{green}，
我们复现了上面的示例。使用分母~$1$ 而非~$2$ 会得到
\showcolorandname{rgb,1:red,1;green,1}，它使用 $1$ 份
\showcolorandname{red} 和 $1$ 份 \showcolorandname{green}。
许多程序允许在 $0,\dotsc,255$ 之间选择分量，因此分母为 $255$ 很有用。
于是，\showcolorandname{rgb,255:red,231;green,84;blue,121} 使用
$231/255$ 红色、$84/255$ 绿色和 $121/255$ 蓝色。这对应于标准
RGB 颜色 $(231,84,121)$。其他示例有
\showcolorandname{rgb,255:red,32;green,127;blue,43}、
\showcolorandname{rgb,255:red,178;green,127;blue,43}、
\showcolorandname{rgb,255:red,169;green,178;blue,43}。

也可以直接使用 RGB 值、HSV 颜色模型、CMY（或 CMYK）模型或
HTML 颜色语法。不过这需要更多编程。我认为这是使用颜色最快
（也可能是最不方便）的方法。例如，
%
\begin{codeexample}[]
\definecolor{color1}{rgb}{1,1,0}
\tikz \fill[color1]
    (0,0) rectangle (1em,0.6em);
\end{codeexample}
%
\noindent 创建具有 $100\%$ \showcolorandname{red}、
$100\%$ \showcolorandname{green} 和 $0\%$ \showcolorandname{blue}
的颜色；
%
\begin{codeexample}[]
\definecolor{color1}{cmyk}{0.6,0.9,0.5,0.1}
\tikz \fill[color1]
    (0,0) rectangle (1em,0.6em);
\end{codeexample}
%
\noindent 创建具有 $60\%$ \showcolorandname{cyan}、
$90\%$ \showcolorandname{magenta}、$50\%$ \showcolorandname{yellow}
和 $10\%$ \showcolorandname{black} 的颜色；

\begin{codeexample}[]
\definecolor{color1}{HTML}{D0B22B}
\tikz \fill[color1]
    (0,0) rectangle (1em,0.6em);
\end{codeexample}
%
\noindent 创建具有 $208/255$ 份红色、$178/255$ 份绿色和 $43$ 份蓝色的
颜色，以标准 HTML 记号指定。更多详情和颜色模型请参阅
|xcolor| 手册~\cite{xcolor}。

|xcolor| 宏包还提供更多组合颜色的方法，其中包括前缀 `|-|'（减号），
它将颜色变为其补色（\showcolorandname{-black}、
\showcolorandname{-white}、\showcolorandname{-red}）或色轮计算。
请参阅 |xcolor| 手册~\cite{xcolor}。
}%

\begin{keylist}{
    /tikz/color=\marg{a color},
    /tikz/draw=\marg{stroke color},
    /tikz/fill=\marg{fill color}%
}
    这些键（通常）用于设置颜色。使用 |color| 同时设置绘制和填充颜色。
    可以直接写 \meta{color name}，而不必写 |color=|\marg{color name}。
    |draw| 和 |fill| 键分别仅设置描边和填充颜色。

    使用 |draw=none| 禁用绘制，使用 |fill=none| 禁用填充。\footnote{到目前为止，
    绘图标记始终具有描边颜色（有些还有填充颜色）。此限制可能在未来
    版本中解除。}
%    This does also work for markers.
%--------------------------------------------------
%\ begin{codeexample}[]
%\begin{tikzpicture}
%\begin{axis}
%    \addplot+[only marks,mark=square*,
%        mark options={fill=red!50!white,draw=none}]
%        {4*x^2 - 2*x +4 };
%\end{axis}
%\end{tikzpicture}
%\end{codeexample}
%--------------------------------------------------

    由于这些键属于 \Tikz{}，完整文档可在 \Tikz{} 手册
    ~\cite[Section ``Specifying a Color'']{tikz} 中找到。
\end{keylist}


\subsubsection{颜色空间}

由于 \PGFPlots{} 依赖 |xcolor|，|xcolor| 中定义颜色空间的所有机制
在此同样适用。

最有用的方法之一是全局颜色空间转换：如果你希望文档仅使用
|cmyk| 颜色空间中的颜色，可以写

\begin{codeexample}[code only]
\usepackage[cmyk]{xcolor}
\usepackage{pgfplots}
\end{codeexample}
%
\noindent 以将整个文档的所有颜色（包括所有底纹）转换为 |cmyk|。

同样的效果可以通过 |xcolor| 语句
\declareandlabel{\selectcolormodel} 实现。
%
\begin{codeexample}[code only]
\selectcolormodel{cmyk}
\end{codeexample}


\subsection{颜色映射}
\label{sec:pgfplots:colormap}

``颜色映射''是一系列颜色，颜色之间平滑过渡。颜色映射常用于在图中
可视化``颜色数据''：此时一个图具有位置坐标 $(x,y)$ 和某个附加标量值
（|point meta|），可用作``颜色数据''。遇到的最小 |point meta| 映射到
|colormap| 的第一个颜色，遇到的最大 |point meta| 值映射到
|colormap| 的最后一个颜色，中间进行插值。

\begin{pgfplotskey}{colormap name=\marg{color map name} (initially hot)}
    将当前颜色映射更改为已定义的、名为 \meta{color map name} 的映射。
    预定义的颜色映射有

    \begin{tabular}{>{\ttfamily}ll}
        viridis & \pgfplotsshowcolormap{viridis} \\
        hot     & \pgfplotsshowcolormap{hot}     \\
    \end{tabular}

    其定义分别见 |colormap/hot| 和 |colormap/viridis| 的文档。
    这些以及更多颜色映射在下文描述。

    颜色映射可用于例如 |scatter| 图和 |surf|ace 图。

    可以使用 |colormap| 创建新的颜色映射（见下文）。
\end{pgfplotskey}

\begin{pgfplotskey}{colormap=\marg{name}\marg{color specification}}
    根据 \meta{color specification} 定义一个名为 \meta{name} 的新颜色映射，
    并使用 |colormap name=|\marg{name} 激活它。

    一个简单的 \meta{color specification} 只是一系列类型为

    \meta{color type}|=(|\meta{color value}|)|

    的颜色定义，以空白符、分号或逗号分隔：
    %
\begin{codeexample}[code only]
\pgfplotsset{colormap={CM}{rgb=(0,0,1) color=(black) rgb255=(238,140,238)}}
\end{codeexample}
\pgfplotsshowcolormapexample{rgb=(0,0,1) color=(black) rgb255=(238,140,238)}

    此处，三个输入颜色分别构成区间的左端、中点和右端。有多种
    \meta{color type} 可用，\meta{color value} 取决于实际的
    \meta{color type}，下文将详细说明。大多数 |colormap| 定义使用
    简单形式，仅列出合适的颜色定义。

    更高级的 \meta{color specification} 同时定义颜色和位置，
    以确定每种颜色的位置。此时，\PGFPlots{} 提供语法

    \meta{color type}|(|\meta{offset}|)=(|\meta{color value}|)|：

\begin{codeexample}[code only]
\pgfplotsset{colormap={CM}{rgb(-500)=(0,0,1) color(0)=(black) rgb255(1500)=(238,140,238)}}
\end{codeexample}
\pgfplotsshowcolormapexample{rgb(-500)=(0,0,1) color(0)=(black) rgb255(1500)=(238,140,238)}

    该语法允许使用非均匀距离在区间上分布颜色。


    \paragraph{兼容性说明：}

    \PGFPlots{} 至版本 $1.13$（含）仅提供对非均匀颜色映射的初步支持。
    需要写 |compat=1.14| 或更高版本才能使用非均匀颜色映射。%
    \index{Color maps!Non uniform, compatibility constraints}

    位置可以是任意数字（或尺寸），\footnote{注意 \PGFPlots{} 至版本 $1.13$（含）
    仅支持范围 $[0,16300]$。} 但每个新颜色的位置必须大于前一个颜色。
    位置可以省略，此时将根据上下文推断：如果前两个颜色没有位置，
    第一个取位置~$0$，第二个取位置 ``$1\text{cm}$''。其后所有颜色取
    最后遇到的网格宽度。注意，非均匀位置在与
    |colormap access=piecewise const| 配合时会产生实际差异。

    精确的输入格式在下一节描述。


    \subsubsection{颜色映射输入格式参考}
    \label{sec:colormap:input:format}

    \meta{color specification} 中的每个条目具有形式

    \meta{color model}|(|\meta{position}|)=(|\meta{arguments}|)| 或

    \meta{special mode}|(|\meta{position}|)=(|\meta{argument}| of |\meta{colormap name}|)|

    其中最常见的形式是使用 \meta{color model} 指定颜色，如
    |rgb=(0,0.5,1)|。\meta{special mode} 在第~\ref{sec:pgfplots:colormaps:based:on:others}
    节（第~\pageref{sec:pgfplots:colormaps:based:on:others} 页）中讨论；
    它们用于访问现有颜色映射的颜色。\meta{position} 参数是
    可选参数，第一个颜色默认为 $0$，第二个颜色默认为 $1\text{cm}$，
    其余颜色使用自动推导的网格间距。\meta{position} 的取值范围是任意的。
    请注意，\PGFPlots{} 仅记住 \meta{position} 的相对距离，而不是
    它们的绝对值。因此，位置为 $0,1,2$ 的颜色映射
    等价于 $0,10,20$、$-10,0,10$ 或 $-1100,-1000,-900$ 的颜色映射。
    \PGFPlots{} 在内部将输入位置映射到 $[0,1000]$ 的范围，并
    使用这些数值。每个新颜色的 \meta{position} 必须
    大于前一个颜色。
    is larger than the preceding one.

    \meta{color model} 的可用选择如下：
    %
    % ATTENTION : a VERY similar description (actually: the same) is in sec:surf:explicit:color
    % maintain both; the syntax is different
    \begin{description}
        \item \declareandlabel{rgb} 接受形式为
            |(|\meta{red}|,|\meta{green}|,|\meta{blue}|)| 的
            \meta{arguments}，每个分量在区间 $[0,1]$ 内，
        \item \declareandlabel{rgb255} 与 |rgb| 类似，但每个分量
            应在区间 $|[0,255]|$ 内，
        \item \declareandlabel{gray} 此时 \meta{arguments} 是区间
            $[0,1]$ 内的单个数字，
        \item \declareandlabel{color} 此时 \meta{arguments} 包含
            预定义（命名）颜色如 `|red|' 或颜色表达式如 `|red!50|'，
        \item \declareandlabel{cmyk} 接受形式为
            |(|\meta{cyan}|,|\meta{magenta}|,|\meta{yellow}|,|\meta{black}|)|
            的 \meta{arguments}，每个分量在区间 $[0,1]$ 内，
        \item \declareandlabel{cmyk255} 与 |cmyk| 相同，但分量应在
            区间 $[0,255]$ 内，
        \item \declareandlabel{cmy} 接受形式为
            |(|\meta{cyan}|,|\meta{magenta}|,|\meta{yellow}|)| 的
            \meta{arguments}，每个分量在区间 $[0,1]$ 内，
        \item \declareandlabel{hsb} 接受形式为
            |(|\meta{hue}|,|\meta{saturation}|,|\meta{brightness}|)| 的
            \meta{arguments}，每个分量在区间 $[0,1]$ 内，
        \item \declareandlabel{Hsb} 与 |hsb| 相同，但 \meta{hue}
            接受区间 $[0,360]$（度）内的值，
        \item \declareandlabel{HTML} 与 |rgb255| 类似，但每个分量
            应为 |00| 到 |FF| 之间的十六进制数，
        \item \declareandlabel{wave} 接受单个波长作为数值参数，
            范围为 $[363,814]$。
    \end{description}

\begin{codeexample}[]
\begin{tikzpicture}
\begin{axis}[
    colormap={bw}{
        gray(0cm)=(0);
        gray(1cm)=(1);
    },
]
    \addplot+ [
        scatter,
        only marks,
        domain=0:8,
        samples=100,
    ] {exp(x)};
\end{axis}
\end{tikzpicture}
\end{codeexample}

    \meta{positions} 的选择会影响处理时间：\meta{positions} 之间均匀
    距离比非均匀距离查找更高效。如果视觉上没有差异，且对数据点数量
    没有影响，优先使用均匀距离的颜色映射而非非均匀映射。注意，
    非均匀映射在与 |colormap access=piecewise constant| 配合时会产生
    巨大差异。\PGFPlots{} 提供了一种简单方法将非均匀颜色定义映射为
    均匀颜色定义：将目标网格宽度写在说明的\emph{第一项}。只要遇到的
    所有 \meta{positions} 都能映射到目标网格，\PGFPlots{} 会自动执行
    插值：\index{Color maps!Non uniform}

\begin{codeexample}[code only]
% non-uniform spacing example: the mesh width is provided as first
% part of the specification.
\pgfplotsset{colormap={violetnew}
    {[1cm] rgb255(0cm)=(25,25,122) color(1cm)=(white) rgb255(5cm)=(238,140,238)}}
\end{codeexample}
\pgfplotsshowcolormapexample{[1cm]
    rgb255(0cm)=(25,25,122)
    color(1cm)=(white)
    rgb255(5cm)=(238,140,238)
}

    \noindent 在最后这个示例中，网格宽度被显式给出，\PGFPlots{} 自行
    插值缺失的网格点。如果所给位置不是网格宽度的整数倍，则会出错。


    \subsubsection{颜色映射的颜色空间}

    \emph{注意：} 如果你已配置 |xcolor| 宏包全局覆盖颜色空间
    （例如在加载 \PGFPlots{} 之前使用 |\usepackage[cmyk]{xcolor}|），
    本节基本多余，见本小节末尾。

    尽管 |colormap| 在\emph{输入}时接受许多颜色空间（事实上接受 |xcolor|
    提供的大多数或全部），但颜色空间的\emph{输出}颜色有严格限制。
    输出颜色空间是 \PGFPlots{} 在两种其他颜色之间进行插值的颜色空间。
    为此，它将输入颜色转换到输出颜色空间。输出颜色空间也称为
    ``颜色映射的颜色空间''。

    |colormap| 支持三种颜色空间：GRAY、RGB 和 CMYK 颜色空间。每次访问
    |colormap| 都需要在其颜色空间中进行线性插值。颜色空间会产生差异：
    不同颜色空间中的颜色可能表示不同，取决于输出设备。许多打印机使用
    CMYK 进行彩色打印，因此提供 CMYK 颜色可能提高彩色打印机上的打印质量。
    RGB 颜色空间常用于显示设备。\PGFPlots{} 中预定义的 |colormap| 都使用 RGB。

    每当创建新的 |colormap|，\PGFPlots{} 会确定一个关联的颜色空间。
    然后，该特定 |colormap| 中的每种颜色都将以其关联的颜色空间表示
    （必要时自动转换颜色）。此外，对 |colormap| 的每次访问都将在其
    关联颜色空间中进行，每次返回的 |mapped color| 也将根据该颜色空间
    表示。此外，由 |shader=interp| 生成的每种底纹都将根据 |colormap|
    的关联颜色空间表示。

    颜色空间选择如下：若 |colormap default colorspace=auto|（初始配置），
    颜色空间取决于 \meta{color specification} 中\emph{第一个}遇到的颜色。
    对于 |rgb| 或 |gray| 或 |color|，关联颜色空间为 RGB
    （与 \PGFPlots{} 早期版本相同）。对于 |cmyk|，关联颜色空间为 CMYK。
    若 |colormap default colorspace| 为 |gray|、|rgb| 或 |cmyk| 之一，
    则使用该特定颜色空间，并自动转换每种颜色。

    \begin{pgfplotskey}{colormap default colorspace=\mchoice{auto,gray,rgb,cmyk} (initially auto)}
        允许设置每个\emph{新建} |colormap| 的颜色空间。
        各选项在上一段中说明。

        无法更改现有 |colormap| 的颜色空间；如需转换，请重新创建。

        若有疑问，宏 \declareandlabel{\pgfplotscolormapgetcolorspace}\marg{name}
        会将 |\pgfplotsretval| 定义为包含现有 |colormap name| 的颜色空间。

        注意，如果你已告知 |xcolor| 全局覆盖颜色空间，该选项\emph{无效}。
        更准确地说，使用
\begin{codeexample}[code only]
\usepackage[cmyk]{xcolor}
\end{codeexample}
        %
        或，替代地，
        %
\begin{codeexample}[code only]
\selectcolormodel{cmyk}
\end{codeexample}
        %
        会导致所有颜色转换为 |cmyk|，\PGFPlots{} 尊重此配置。
        因此，这两条语句都会使所有颜色在所需颜色空间中插值，
        所有输出颜色都将使用该颜色空间。这通常正是你所需要的。
    \end{pgfplotskey}


    \subsubsection{预定义颜色映射}

    可用的颜色映射如下所示。
\end{pgfplotskey}

\begin{stylekey}{/pgfplots/colormap/viridis}
    该样式安装由 Stèfan van der Walt 和 Nathaniel Smith 为
    |Matplotlib| 定义的 ``viridis'' 颜色映射。它被指定为
    |Matplotlib| 从版本~2.0 起的默认 |colormap|，并以
    CC0\footnote{\url{https://creativecommons.org/about/cc0}} 发布。

    |viridis| 选项是 \PGFPlots{} 中包含的降采样副本。

\begin{codeexample}[code only]
\pgfplotsset{
    colormap name=viridis,
}
\end{codeexample}
    \pgfplotsshowcolormap{viridis}

    与其他选项相比，该 |colormap| 具有明显更好的特性：
    %
    \begin{itemize}
        \item 其颜色分布在感知上是均匀的（参见下方链接中的定义），
        \item 它适用于中等程度的色盲，
        \item 黑白打印时仍然良好。
    \end{itemize}
    %
    有关这些特性的详情，请参阅
    \url{http://bids.github.io/colormap}。

    请使用 |colormap name=viridis| 选项，因为它使用预定义颜色映射，
    而 |colormap/viridis| 会重新定义它。

    |colormaps| 库中还有一个 |viridis| 的高分辨率副本，
    称为 |colormap/viridis high res|。它接近作者原始分辨率，
    但视觉上与 |viridis| 几乎相同，且在 \TeX{} 中占用更少资源。
\end{stylekey}

\begin{stylekey}{/pgfplots/colormap/hot}
    该样式安装颜色映射
    %
\begin{codeexample}[code only]
\pgfplotsset{
    colormap={hot}{color(0cm)=(blue); color(1cm)=(yellow); color(2cm)=(orange); color(3cm)=(red)}
}
\end{codeexample}

    \pgfplotsshowcolormap{hot}

    这是一个预配置的颜色映射。
    %
\end{stylekey}

\begin{stylekey}{/pgfplots/colormap/hot2}
    该样式等价于
    %
\begin{codeexample}[code only]
\pgfplotsset{
    /pgfplots/colormap={hot2}{[1cm]rgb255(0cm)=(0,0,0) rgb255(3cm)=(255,0,0)
        rgb255(6cm)=(255,255,0) rgb255(8cm)=(255,255,255)}
}
\end{codeexample}

    \pgfplotsshowcolormap{hot2}

    注意，该特定选项随 \PGFPlots{} 直接提供，无需为该值加载 |colormaps| 库。

    \matlabcolormaptext
\end{stylekey}

\begin{stylekey}{/pgfplots/colormap/jet}
    该样式等价于
    %
\begin{codeexample}[code only]
\pgfplotsset{
    /pgfplots/colormap={jet}{rgb255(0cm)=(0,0,128) rgb255(1cm)=(0,0,255)
        rgb255(3cm)=(0,255,255) rgb255(5cm)=(255,255,0) rgb255(7cm)=(255,0,0) rgb255(8cm)=(128,0,0)}
}
\end{codeexample}

    \pgfplotsshowcolormap{jet}

    \matlabcolormaptext
\end{stylekey}

\begin{stylekey}{/pgfplots/colormap/blackwhite}
    该样式等价于
    %
\begin{codeexample}[code only]
\pgfplotsset{
    colormap={blackwhite}{gray(0cm)=(0); gray(1cm)=(1)}
}
\end{codeexample}

    \pgfplotsshowcolormap{blackwhite}
\end{stylekey}

\begin{stylekey}{/pgfplots/colormap/bluered}
    该样式等价于
    %
\begin{codeexample}[code only]
\pgfplotsset{
    colormap={bluered}{
        rgb255(0cm)=(0,0,180); rgb255(1cm)=(0,255,255); rgb255(2cm)=(100,255,0);
        rgb255(3cm)=(255,255,0); rgb255(4cm)=(255,0,0); rgb255(5cm)=(128,0,0)}
}
\end{codeexample}

    \pgfplotsshowcolormap{bluered}

\begin{codeexample}[]
\begin{tikzpicture}
\begin{axis}[colormap/bluered]
    \addplot+ [scatter,
         scatter src=x,samples=50]
        {sin(deg(x))};
\end{axis}
\end{tikzpicture}
\end{codeexample}

    \paragraph{说明：}

    样式 |bluered|（重新）定义颜色映射并激活它。如果在导言区调用
    |\pgfplotsset{colormap/bluered}|（仅创建一次颜色映射），并在需要时
    使用 |colormap name=bluered|，\TeX{} 会稍快。本说明适用于后续
    每个颜色映射样式。但你可以直接忽略本说明。
\end{stylekey}

\begin{stylekey}{/pgfplots/colormap/cool}
    该样式等价于
    %
\begin{codeexample}[code only]
\pgfplotsset{
    colormap={cool}{rgb255(0cm)=(255,255,255); rgb255(1cm)=(0,128,255); rgb255(2cm)=(255,0,255)}
}
\end{codeexample}

    \pgfplotsshowcolormap{cool}
\end{stylekey}

\begin{stylekey}{/pgfplots/colormap/greenyellow}
    该样式等价于
    %
\begin{codeexample}[code only]
\pgfplotsset{
    colormap={greenyellow}{rgb255(0cm)=(0,128,0); rgb255(1cm)=(255,255,0)}
}
\end{codeexample}

    \pgfplotsshowcolormap{greenyellow}
\end{stylekey}

\begin{stylekey}{/pgfplots/colormap/redyellow}
    该样式等价于
    %
\begin{codeexample}[code only]
\pgfplotsset{
    colormap={redyellow}{rgb255(0cm)=(255,0,0); rgb255(1cm)=(255,255,0)}
}
\end{codeexample}

    \pgfplotsshowcolormap{redyellow}
\end{stylekey}

\begin{stylekey}{/pgfplots/colormap/violet}
    该样式等价于
    %
\begin{codeexample}[code only]
\pgfplotsset{
    colormap={violet}{rgb255=(25,25,122) color=(white) rgb255=(238,140,238)}
}
\end{codeexample}

    \pgfplotsshowcolormap{violet}
\end{stylekey}

\begin{command}{\pgfplotscolormaptoshadingspec\marg{colormap name}\marg{right end size}\marg{\textbackslash macro}}
    该命令将颜色映射转换为 \PGF{} 底纹的颜色说明。它可用于
    |\pgfdeclare...shading| 等命令（详情见 \PGF{} 手册~\cite{tikz}）。

    第一个参数是（已定义的）颜色映射名称，第二个是说明最右端的尺寸。
    结果将存储在 \meta{\textbackslash macro} 中。
    %
\begin{codeexample}[]
    % convert `hot' -> \result
    \pgfplotscolormaptoshadingspec{hot}{8cm}\result
    % define and use a shading in pgf:
    \def\tempb{\pgfdeclarehorizontalshading{tempshading}{1cm}}%
    % where `\result' is inserted as last argument:
    \expandafter\tempb\expandafter{\result}%
    \pgfuseshading{tempshading}%
\end{codeexample}
    %
    结果 \meta{\textbackslash macro} 的使用有点底层。


    \paragraph{注意：}

    \PGF{} 底纹始终根据 RGB 颜色空间表示。因此，使用此方法时，
    即使是 CMYK \meta{colormap name} 也会产生 RGB 底纹说明。\footnote{
    如果 \PGF{} 将来支持 CMYK 底纹，而你仍看到本说明，可以在导言区添加
    宏定义 \texttt{\textbackslash def\textbackslash
    pgfplotscolormaptoshadingspectorgb\{0\}}。}
\end{command}

\begin{command}{\pgfplotscolorbardrawstandalone\oarg{options}}
    该命令绘制一个 |tikzpicture| 和其中使用当前颜色条设置的
    |colorbar|。其目的是简化文档。

    由于此 |colorbar| 是``独立''图片，它定义以下选项
    %
\begin{codeexample}[code only]
    point meta min=0,
    point meta max=1000,
    parent axis width/.initial=6cm,
    parent axis height/.initial=6cm,
\end{codeexample}
    %
    然后才求值 \meta{options} 并绘制颜色条。
\end{command}

\begin{command}{\pgfplotscolormaptodatafile\oarg{options}\marg{colormap name}\marg{output file}}
    允许将 |colormap| 数据导出到文件。
    %
\begin{codeexample}[]
\begin{tikzpicture}
\begin{axis}[y=1cm,table/col sep=comma]
        \pgfplotscolormaptodatafile{hot}{hot.dat}
    \addplot[red,mark=|]   table[y index=1] {hot.dat};
    \addplot[green,mark=|] table[y index=2] {hot.dat};
    \addplot[blue,mark=|]  table[y index=3] {hot.dat};
\end{axis}
\end{tikzpicture}
\end{codeexample}

\begin{codeexample}[]
\begin{tikzpicture}
\begin{axis}[y=1cm,table/col sep=comma]
    \pgfplotscolormaptodatafile{viridis}{viridis.dat}
    \addplot[red]   table[y index=1] {viridis.dat};
    \addplot[green] table[y index=2] {viridis.dat};
    \addplot[blue]  table[y index=3] {viridis.dat};
\end{axis}
\end{tikzpicture}
\end{codeexample}

    有效的 \meta{options} 有
    %
    \begin{pgfplotskey}{/pgfplots/colormap/output each nth=\meta{num} (initially 1)}
        允许通过仅写入每第 \meta{num} 个条目来降采样颜色映射。
    \end{pgfplotskey}

    \begin{pgfplotskey}{/pgfplots/colormap/output format=\mchoice{cvs,native} (initially csv)}
        选项 \declaretext{csv} 生成 CSV 文件，其中第一列是颜色映射的偏移，
        其余列是颜色分量。

        选项 \declaretext{native} 生成 \TeX{} 文件，可 |\input| 以定义
        颜色映射供 \PGFPlots{} 使用。
    \end{pgfplotskey}
\end{command}

注意，相关库中\emph{有更多可用颜色映射选择}，即 |colorbrewer| 库和
|colormaps| 库，需通过 |\usepgfplotslibrary{colorbrewer,colormaps}| 加载。


\subsubsection{基于其他颜色映射构建颜色映射}
\label{sec:pgfplots:colormaps:based:on:others}

{
\tikzset{
    /pdflinks/search key prefixes in/.add={}{,/pgfplots/of colormap/},
}
形如 |colormap=|\marg{name}\marg{color specification} 的颜色映射定义
通常由单个颜色组成的 \meta{color specifications} 构成，
每个颜色有自己的 \meta{color model}。如第~\ref{sec:colormap:input:format}
节所述，\meta{color specification} 中每个条目具有形式

    \meta{color model}|(|\meta{position}|)=(|\meta{arguments}|)| 或

    \meta{special mode}|(|\meta{position}|)=(|\meta{argument}| of |\meta{colormap name}|)|。

\noindent 本节说明如何利用 \meta{special mode} 基于现有 |colormap|
构建 |colormap|。为此，\PGFPlots{} 在 \meta{color specification} 内
提供以下值：

\begin{enumerate}
    \item
        \declareandlabel{samples of colormap}|(|\meta{position}|)=(|\meta{number}| of |\meta{colormap name}|)| 或\\
        \declaretext{samples of colormap}|(|\meta{position}|)=(|\meta{number}|)| 或\\
        \declaretext{samples of colormap}|(|\meta{position}|)={|\meta{number}| of |\meta{colormap name}|, |\meta{options}|}|%
            \index{samples of colormap!In colormap specification}%

        该方法接受 \meta{colormap name} 作为输入，从中采样 \meta{number}
        个颜色（使用均匀网格宽度），并插入当前正在构建的 |colormap|。
        它等价于选项 |colors of colormap||(|\meta{position}|)=(0,h,...,1000)|，
        其中 |h| 的选择使你得到位于等距网格上的 \meta{number} 个样本。

\begin{codeexample}[]
\pgfplotscolorbardrawstandalone[
    colormap={example}{
        samples of colormap=(4 of viridis)
    },
    colorbar horizontal,
    colormap access=const,
]
\end{codeexample}

\begin{codeexample}[]
\pgfplotscolorbardrawstandalone[
    colormap={example}{
        samples of colormap=(4 of viridis)
    },
    colorbar horizontal,
    colormap access=map,
]
\end{codeexample}

        特殊后缀 ``| of |\meta{colormap name}'' 是可选的；
        它默认为 |colormap name| 的当前值：
        %
\begin{codeexample}[]
\pgfplotscolorbardrawstandalone[
    colormap={example}{
        samples of colormap=(4)
    },
    colorbar horizontal,
    colormap access=const,
]
\end{codeexample}

\begin{codeexample}[]
\pgfplotscolorbardrawstandalone[
    colormap={example}{
        samples of colormap=(4)
    },
    colorbar horizontal,
    colormap access=map,
]
\end{codeexample}

        参数可以用圆括号或花括号包围，两者都接受。\PGFPlots{} 也接受
        参数内的一系列选项；如果需要选项，最好使用花括号：
        %
\begin{codeexample}[]
\pgfplotscolorbardrawstandalone[
    colormap={example}{
        samples of colormap={
            5 of viridis,
            target pos={0,800,850,950,1000},
        }
    },
    colorbar horizontal,
    colormap access=map,
]
\end{codeexample}

\begin{codeexample}[]
\pgfplotscolorbardrawstandalone[
    colormap={example}{
        % simpler alternative:
        % of colormap={viridis,target pos={...}}
        samples of colormap={
            5 of viridis,
            target pos={0,400,500,700,800,1000},
            sample for=const,
        }
    },
    colorbar horizontal,
    colormap access=const,
]
\end{codeexample}

        参数内选项的使用在下一小节讨论，见 |of colormap| 与 |target pos|。

        与普通颜色定义一样，\meta{position} 参数是可选的，可以省略。
        如果给出，则用于列表中第一个遇到的条目，其余自动推断。
        如果同时给出 |target pos| 和 \meta{position}，则忽略
        \meta{position}。

        注意，\PGFPlots{} 为 |target pos| 和采样颜色映射提供了特殊语法：
        %
\begin{codeexample}[]
\pgfplotscolorbardrawstandalone[
    colormap={example}{
        of colormap={
            viridis,
            target pos={0,400,500,700,800,1000},
            sample for=const,
        }
    },
    colorbar horizontal,
    colormap access=const,
]
\end{codeexample}
        %
        该语法同样从源颜色映射（此处为 |viridis|）采样颜色。
        它选择足够的样本以满足给定的 |target pos|；
        详情见 `|of colormap|' 的文档。
    \item
        \declareandlabel{index of colormap}|(|\meta{position}|)=(|\meta{index}| of |\meta{colormap name}|)| 或\\
        \declaretext{index of colormap}|(|\meta{position}|)=(|\meta{index}|)|%
            \index{index of colormap!In colormap specification}%

        该键允许标识 \meta{colormap name} 中的\emph{单个}颜色，
        并将其用作 \meta{color specification} 的一部分。
        \meta{index} 是数值索引 $0,1,\dotsc,N-1$，其中 $N$ 是
        \meta{colormap name} 的大小。

\begin{codeexample}[]
\pgfplotscolorbardrawstandalone[
    colormap={example}{
        color=(green),
        index of colormap=(2 of viridis)
    },
    colorbar horizontal,
    colormap access=const,
]
\end{codeexample}

        注意，|index of colormap| 也可用作绘图操作的键：

        \begin{pgfplotskey}{index of colormap=\meta{index} of \meta{colormap name}}
            使用特殊颜色名 `|.|' 选择 \meta{colormap name} 的指定颜色，
            并赋值 |color=.|：
            %
\begin{codeexample}[]
\pgfplotsset{
 colormap={example}{color=(red) color=(blue)}}
\tikz \draw[index of colormap=0 of example] (0,0) -- (1,0);
\end{codeexample}
            %
\begin{codeexample}[]
\pgfplotsset{
 colormap={example}{color=(red) color=(blue)}}
\tikz \fill[index of colormap=1 of example] (0,0) rectangle (1,1);
\end{codeexample}
            %
            注意，某些绘图指令需要的不只是设置 |color|。此时，
            你必须在该键之后引用颜色 `|.|'：
            %
\begin{codeexample}[]
\pgfplotsset{
 colormap={example}{color=(red) color=(blue)}}
\tikz \shade[
    index of colormap=0 of example,left color=.,
    color of colormap=500 of example,right color=.] (0,0) rectangle (1,1);
\end{codeexample}
        \end{pgfplotskey}

        |samples of colormap| 的所有特殊说明（如花括号、选项列表支持、
        位置）在此同样适用。
    \item
        \declareandlabel{indices of colormap}|(|\meta{position}|)=(|\meta{list of indices}| of |\meta{colormap name}|)| 或\\
        \declaretext{indices of colormap}|(|\meta{position}|)=(|\meta{list of indices}|)|%
            \index{indices of colormap!In colormap specification}%

        一个便捷键，等价于一系列 |index of colormap||(|\meta{position}|)|，
        \meta{list} 中每个元素一个。\meta{list} 使用 |\foreach| 求值。
        注意，参数需要用圆括号包围。

\begin{codeexample}[]
\pgfplotscolorbardrawstandalone[
    colormap={example}{
        indices of colormap=(0,5,10,12,
          \pgfplotscolormaplastindexof{viridis}
          of viridis)
    },
    colorbar horizontal,
    colormap access=const,
]
\end{codeexample}

        |samples of colormap| 的所有特殊说明（如花括号、选项列表支持、
        位置）在此同样适用。
    \item
        \declareandlabel{color of colormap}|(|\meta{position}|)=(|\meta{value}| of |\meta{colormap name}|)| 或\\
        \declaretext{color of colormap}|(|\meta{position}|)=(|\meta{value}|)|%
            \index{color of colormap!In colormap specification}%

        该键允许在 \meta{colormap name} 内插值一个颜色，并将结果用作
        \meta{color specification} 的一部分。插值点是范围 $[0,1000]$
        内的浮点数，使用 |colormap access=map| 插值
        （即使用分段线性插值）。

        注意，|color of colormap| 也可用作绘图操作的键：

        \begin{pgfplotskey}{color of colormap=\meta{value} of \meta{colormap name}}
            使用特殊颜色名 `|.|' 选择 \meta{colormap name} 的指定颜色，
            并赋值 |color=.|：
            %
\begin{codeexample}[]
\pgfplotsset{
 colormap={example}{color=(red) color=(blue)}}
\tikz \shade[
    color of colormap=250 of example,left color=.,
    color of colormap=500 of example,right color=.] (0,0) rectangle (1,1);
\end{codeexample}
        \end{pgfplotskey}

        |samples of colormap| 的所有特殊说明（如花括号、选项列表支持、
        位置）在此同样适用。
    \item
        \declareandlabel{colors of colormap}|(|\meta{position}|)=(|\meta{list}| of |\meta{colormap name}|)| 或\\
        \declaretext{colors of colormap}|(|\meta{position}|)=(|\meta{list}|)| 或\\
        \declaretext{colors of colormap}|(|\meta{position}|)={|\meta{list}| of |\meta{colormap name}|, |\meta{options}|}|%
            \index{colors of colormap!In colormap specification}%

        一个便捷键，等价于一系列 |color of colormap||(|\meta{position}|)|，
        \meta{list} 中每个元素一个。\meta{list} 使用 |\foreach| 求值。
        注意，参数需要用圆括号包围。

\begin{codeexample}[]
\pgfplotscolorbardrawstandalone[
    colormap={example}{
      colors of colormap=(0,400,800,900,
          1000 of viridis)
    },
    colorbar horizontal,
    colormap access=const,
]
\end{codeexample}

        注意，|colors of colormap| 在 |cycle list| 定义期间也可用。

        |samples of colormap| 的所有特殊说明（如花括号、选项列表支持、
        位置）在此同样适用。
    \item \declaretext{of colormap}|(|\meta{position}|)=(|\meta{colormap name}|, |\meta{options}|)|

        通过从输入 |colormap name| 中\emph{采样}足够颜色来构建新的
        |colormap|。该模式是从另一个 |colormap| 采样颜色的两种方法之一；
        如果你想分配 |target pos|，应优先使用它。另见
        |samples of colormap|；如果你只想指定样本数量，它更简单。

        语法 `|of colormap=|' 最适合想从另一个颜色映射
        \emph{在特定位置} 采样颜色时：
        %
\begin{codeexample}[]
\pgfplotscolorbardrawstandalone[
    colormap={example}{
        of colormap={
            viridis,
            target pos={0,400,500,700,800,1000},
        },
    },
    colorbar horizontal,
    colormap access=map,
]
\end{codeexample}
        %
        此处，\PGFPlots{} 解析 \meta{options}。每个未识别选项要么被视为
        `|colormap name|' 的值，要么被视为定义 |colormap|（如
        |colormap/PuOr|）的样式参数。前一个示例识别出
        |colormap name=viridis|（因为 |viridis| 不是已知选项）并分配
        |target pos|。颜色映射由 $6$ 个样本构成，这些样本取自 |viridis|。
        注意，在此上下文中无需指定 \meta{colormap name}。如果缺失，
        将使用 |colormap name| 的当前值（即当前颜色映射）。

        类似地，`|of colormap|' 可以采样 |colormap| 以配合
        |colormap access=const| 使用：
        %
\begin{codeexample}[]
\pgfplotscolorbardrawstandalone[
    colormap={example}{
        of colormap={
            viridis,
            target pos={0,400,500,700,800,1000},
            sample for=const,
        },
    },
    colorbar horizontal,
    colormap access=const]
\end{codeexample}
        %
        注意，在此上下文中\emph{需要}写 |sample for=const|，
        以便 \PGFPlots{} 知道结果要配合 |colormap access=const| 使用。
        详情见下文。

        与所有颜色映射构建块一样，|target pos| 可以有任意数字范围。
        不过，绝对值没有意义；它们始终被映射到范围 $[0,1000]$。
        就 |colormap| 本身而言，只有它们的相对距离重要：
        %
\begin{codeexample}[]
\pgfplotscolorbardrawstandalone[
    colormap={example}{
        of colormap={viridis,
            target pos={-10,-3,0,1,2},
            sample for=const,
        },
    },
    colorbar horizontal,
    colormap access=const]
\end{codeexample}
        %
        但是，如果 |point meta| 范围等于颜色映射定义范围，
        你会看到它们完全吻合：
        %
\begin{codeexample}[]
\pgfplotscolorbardrawstandalone[
    colormap={example}{
        of colormap={viridis,
            target pos={-10,-3,0,1,2},
            sample for=const,
        },
    },
    point meta min=-10,point meta max=2,
    colorbar horizontal,
    colormap access=const]
\end{codeexample}
        %
        注意，你可以在 |colorbar| 样式内使用 |xtick=data| 或 |ytick=data|，
        以将刻度标签放在 |colormap| 的位置处：
        %
\begin{codeexample}[]
\pgfplotscolorbardrawstandalone[
    colormap={example}{
        of colormap={viridis,
            target pos={-10,-3,0,1,2},
            sample for=const,
        },
    },
    point meta min=-10,point meta max=2,
    colorbar horizontal,
    colorbar style={xtick=data},
    colormap access=const]
\end{codeexample}

        最后，`|of colormap|' 可用于\emph{复制}另一个 |colormap|：
        如果没有关于如何构建新 |colormap| 的合适提示，
        \PGFPlots{} 复制输入映射：
\begin{codeexample}[]
\pgfplotscolorbardrawstandalone[
    colormap={example}{
        of colormap={viridis},
    },
    colorbar horizontal,
    colormap access=const]
\end{codeexample}

        \meta{options} 被所有颜色映射构建块接受：如果需要，你也可以在
        |samples of colormap| 或 |colors of colormap| 之后指定它们。

        以下选项可用：

        \begin{pgfplotskey}{of colormap/target pos=\marg{position(s)} (initially empty)}
            允许为零个、一个或多个相关颜色定义 \meta{position}。
            空值表示在定义中使用圆括号 `|(|\meta{position}|)|' 中的参数，
            例如 |color of colormap(100)=(4)| 中的值 $100$。如果只涉及
            一种颜色，这就足够了。但是，如果 |target pos| 有非空值，
            它会覆盖 |(|\meta{position}|)| 参数。

            键 |target pos| 主要用于为多个颜色定义提供位置，
            特别是在
            `|of colormap|'. 在此上下文中，使用所指定数量的 \meta{position(s)}。

            如果你有 |colors of colormap| 和非空的
            |target pos| 值，\PGFPlots{} 将使用尽可能多的可用 |target pos|itions。
            如果数量太少而无法覆盖所有涉及的颜色，则
            剩余的位置将自动推导。更精确地说：
            \PGFPlots{} 将使用当前网格宽度作为到
            前一个位置的距离。``当前网格宽度'' 被定义为
            相邻位置之间的最小差值。多余的位置
            将被忽略（并在日志文件中写入警告）。


            注意，|target pos| 与 |colormap access=const| 一起使用时需要注意：
            此时，有\emph{区间}颜色恒定，位置成为区间边界。因此，
            |target pos| 中 $N$ 个元素构成 $N-1$ 种不同颜色！

            注意，|target pos| 作为 |samples of colormap| 的选项也是可能的。
            但是，你必须指定所有所需位置，数量不匹配则为错误。如果还将
            此情况与 |sample for=const| 组合，$N$ 个样本需要 $N+1$ 个位置。
        \end{pgfplotskey}

        \begin{pgfplotskey}{of colormap/sample for=\mchoice{default,const} (initially default)}
            该键允许配置 |samples of colormap|，使生成的 |colormap|
            适用于 |colormap access| 的特定值。

            特别是，它允许优化依赖\emph{采样}的 \meta{colormap definitions}，
            以适应 |colormap access=const| 的情况。

            注意，|colormap access|$\neq$|const| 的情况通常无需修改此键，
            使用默认值效果最佳。

            策略 |colormap access=const| 意味着颜色与整个\emph{区间}关联，
            而不再与 |colormap| 中的单个位置关联。在此上下文中，
            你必须写 |sample for=const|，以便 \PGFPlots{} 正确处理。

            以下是 |colormap| 参考及其与 |sample for| 的关系：
            %
            \begin{itemize}
                \item 一般来说，|colormap| 可与 |colormap access| 的任何值
                    配合使用。如果 \meta{colormap definition} 中没有位置
                    （既不在圆括号中，也不在 |target pos| 中），
                    则无需特别注意即可直接使用。
                    %
\begin{codeexample}[]
\pgfplotscolorbardrawstandalone[
    colormap={example}{
        samples of colormap={
            5 of viridis,
        },
    },
    colorbar horizontal,
    colormap access=const]
\end{codeexample}
                    %
                \item 一旦有了位置，如果颜色映射在 |colormap access=const|
                    上下文中使用，\PGFPlots{} 就会产生冲突：它无法扩大
                    显示范围，但又无法同时尊重给定位置\emph{和}所需的
                    颜色数量！它只能选择省略一种颜色，或添加一个人工位置
                    并重新缩放所有其他位置。
                \item 在存在 |colormap access=const| 和输入位置时，
                    \PGFPlots{} 的默认策略是：\PGFPlots{} 选择保留输入位置
                    并省略一种颜色。
\begin{codeexample}[]
\pgfplotscolorbardrawstandalone[
    colormap={example}{
        of colormap={
            viridis,
            target pos={0,500,850,950,1000},
        },
    },
    colorbar horizontal,
    colormap access=const]
\end{codeexample}

                    技术上，这意味着它隐式设置
                    |colormap access/extra interval width=0|
                    （默认是 |colormap access/extra interval width=h|）。
                    与 $1000$ 关联的颜色没有区间，实际上不可见。

                    恢复 |colormap access/extra interval width=h| 会导致
                    位置重新缩放，这通常不是你想要的：
                    %
\begin{codeexample}[]
\pgfplotscolorbardrawstandalone[
    colormap={example}{
        samples of colormap={
            5 of viridis,
            target pos={0,500,850,950,1000},
        },
    },
    colormap access/extra interval width=h,
    colorbar horizontal,
    colormap access=const]
\end{codeexample}
                    %
                \item 最后，\PGFPlots{} 提供 |sample for=const|。

                    对 `|of colormap|' 应用 |sample for=const| 可确保
                    尊重任何给定位置。为此，它减少可用颜色数量，
                    但修改采样过程，使整个输入颜色范围在结果中可见：
                    %
\begin{codeexample}[]
\pgfplotscolorbardrawstandalone[
    colormap={example}{
        of colormap={
            viridis,
            target pos={0,500,850,950,1000},
            sample for=const,
        },
    },
    colorbar horizontal,
    colormap access=const]
\end{codeexample}
                    %
                    此处，结果有 $4$ 种颜色，但最右端颜色也是输入映射
                    |viridis| 的最右端颜色。注意，|of colormap| 已为位置
                    $1000$ 分配了一种不可见颜色：该虚拟颜色被定义为
                    |viridis| 的最右端颜色。这意味着，如果你用
                    |sample for=const| 定义 |colormap| 并配合
                    |colormap access=map| 使用，区间右端会比预期长：
                    %
\begin{codeexample}[]
\pgfplotscolorbardrawstandalone[
    colormap={example}{
        of colormap={
            viridis,
            target pos={0,500,850,950,1000},
            sample for=const,
        },
    },
    colorbar horizontal,
    colormap access=map]
\end{codeexample}

                    应用于 |samples of colormap| 的策略 |sample for=const|
                    以类似方式工作。请记住，此组合要求要么省略
                    |target pos|，要么添加\emph{恰好} $N+1$ 个目标位置：
                    %
\begin{codeexample}[]
\pgfplotscolorbardrawstandalone[
    colormap={example}{
        samples of colormap={
            5 of viridis,
            target pos={0,100,500,850,950,1000},
            sample for=const,
        },
    },
    colorbar horizontal,
    colormap access=const]
\end{codeexample}

                    这会正确尊重位置并完全尊重输入颜色映射的范围。
            \end{itemize}

            每当 \meta{colormap definition} 涉及采样时（即仅对
            |of colormap| 和 |samples of colormap|），键 |sample for|
            就会生效。它不适用于显式固定的颜色。其效果是样本数量
            减少~$1$，最后提供的样本被复制一次。

            另见 |colormap access/extra interval width|。
        \end{pgfplotskey}

        \begin{pgfplotskey}{of colormap/colormap access=\marg{argument}}
             |/pgfplots/colormap access=|\marg{argument} 的别名。
        \end{pgfplotskey}

        \begin{pgfplotskey}{of colormap/source range=\marg{min:max} (initially 0:1000)}
             为基于插值的说明（即 |color of colormap| 和
             |const color of colormap|）定义源范围。它默认为 |0:1000|，
             意味着只能提供区间 $[0,1000]$ 内的值。

             更改该值允许使用任意数字范围来标识数字。
             %
\begin{codeexample}[]
\pgfplotscolorbardrawstandalone[
    colormap={example}{
        colors of colormap={
            -30000,-10000,0,100000,110000 of viridis,
            source range=-30000:120000,
        },
    },
    colorbar horizontal,
    colormap access=const]
\end{codeexample}
            %
        \end{pgfplotskey}

%        \begin{pgfplotskey}{of colormap/source name=\marg{colormap name} (initially empty)}
%            The name of the source color map. An empty value means ``use the
%            current value of |colormap name|''.
%        \end{pgfplotskey}

        \noindent 注意，这些构建块可以根据需要任意组合。
        这允许组合不同的颜色映射：
        %
\begin{codeexample}[]
\pgfplotscolorbardrawstandalone[
    point meta min=-7046,
    point meta max=2895,
    colormap={whiteblue}{color=(blue) color=(white)},
    colormap={gb}{color=(green) color=(yellow)
        color=(brown)},
    colormap={CM}{
        of colormap={
            whiteblue,
            target pos={-7046,-6000,-5000,-3000,
                -1000,-750,-500,-250,-100,-50,0},
            sample for=const,
        },
        of colormap={
            gb,
            target pos={10,100,200,500,1000,1100,
                1200,1500,2000,2895},
            sample for=const,
        },
    },
    colorbar horizontal,
    colormap access=const]
\end{codeexample}
        %
        \noindent 前一个示例首先定义了两个简单颜色映射 |whiteblue| 和 |gb|。
        它们仅用作构建块，不用于可视化。最后，颜色映射 |CM| 由两个
        |of colormap| 说明构成，它们对构建块重新采样并分配特定的
        |target pos|itions。

        \begin{pgfplotskeylist}{%
            of colormap/target pos min=\marg{lower limit} (initially empty),
            of colormap/target pos min*=\marg{lower limit} (initially empty),
            of colormap/target pos max=\marg{upper limit} (initially empty),
            of colormap/target pos max*=\marg{upper limit} (initially empty)%
        }
            这些键允许\emph{修改} |target pos| 的参数。
            它们的主要用途是简化样式定义的编写。因此，大多数用户可能
            想忽略这些键并跳过其文档。如果直接使用 |target pos| 的参数，
            这些键就不必要。

            不带星号的版本（|target pos min| 和 |target pos max|）
            丢弃 |target pos| 中所有超出边界的元素。带星号的版本
            （|target pos min*| 和 |target pos max*|）也丢弃所有不匹配项，
            但它们确保过滤后边界成为 |target pos| 的一部分。

            如果你想从先前定义的 |target pos| 中选择匹配位置，
            这些过滤边界会很方便，例如当同一个 |target pos| 是可复用
            样式的一部分时：
            %
\begin{codeexample}[width=8cm]
\pgfplotsset{
    of colormap/ocean height/.style={
        target pos={-12000,-10000,-6000,-5000,-3000,
            -1000,-750,-500,-250,-100,-50,0,10,100,
            200,500,1000,1100,1200,1500,2000,4000,
            6000,8000},
    },
}
\pgfplotscolorbardrawstandalone[
  point meta min=-7046,
  point meta max=2895,
  colormap={whiteblue}{color=(blue) color=(white)},
  colormap={gb}{color=(green) color=(yellow)
      color=(brown)},
  colormap={CM}{
      of colormap={
        whiteblue,
        ocean height,
        target pos min*=
          \pgfkeysvalueof{/pgfplots/point meta min},
        target pos max=0,
        sample for=const,
      },
      of colormap={
        gb,
        ocean height,
        target pos min=0.1,
        target pos max*=
          \pgfkeysvalueof{/pgfplots/point meta max},
        sample for=const,
      },
  },
  colorbar horizontal,
  colormap access=const]
\end{codeexample}
            %
            \noindent 该示例在文档某处定义了一个名为 `|ocean height|' 的
            样式，其中包含合适的位置列表。然后，它定义了颜色映射 `|CM|'，
            该颜色映射使用这些键——但仅在范围 $[-7046,2895]$ 内，
            并以两个其他 |colormap| 的特殊组合。第一个 |of colormap|
            说明仅选择落在范围 $[-7046,0]$ 内的目标位置，并确保
            $-7046$ 实际成为 |target pos| 的元素。第二个 |of colormap|
            说明选择范围 $[0.1,2895]$ 内的所有位置，并确保 $2895$
            成为 |target pos| 的元素。注意，$0.1$ 仅用作指示符，
            表示\emph{不}再次选择 $0$。因此，该选择本质上等价于
            %
\begin{codeexample}[code only]
...
target pos={-7046,-6000,-5000,-3000,-1000,-750,-500,-250,-100,-50,0}
...
target pos={10,100,200,500,1000,1100,1200,1500,2000,2895}
\end{codeexample}
            %
            区别仅在于预定义样式定义了合适的位置列表。从 |whiteblue|
            取的\emph{颜色}与写 |samples of colormap={11}| 时相同，
            即它们从输入颜色映射均匀抽取。只是它们在结果中的目标位置
            被修改。从 |bg| 取的颜色同理；它们被均匀抽取并移到规定边界。

            \PGFPlots{} 使用这些键来实现 |contour filled|。

            注意，颜色映射定义不绑定于特定坐标，尽管这在 |contour|
            图的上下文中很有意义。原则上，你可以在任何上下文中使用
            生成的颜色映射，就像所有其他颜色映射一样：
            %
\begin{codeexample}[]
\pgfplotsset{
  colormap={whiteblue}{color=(blue) color=(white)},
  colormap={gb}{color=(green) color=(yellow)
      color=(brown)},
  of colormap/ocean height/.style={
      target pos={-12000,-10000,-6000,-5000,-3000,
          -1000,-750,-500,-250,-100,-50,0,10,100,
          200,500,1000,1100,1200,1500,2000,4000,
          6000,8000},
  },
  colormap={CM}{
      of colormap={
          whiteblue,
          ocean height,
          target pos max=0,
          sample for=const,
      },
      of colormap={
          gb,
          ocean height,
          target pos min=0.1,
          sample for=const,
      },
  },
}
\begin{tikzpicture}
\begin{axis}[small,colorbar]
    \addplot3[surf] {x*y};
\end{axis}
\end{tikzpicture}
~
\begin{tikzpicture}
\begin{axis}[small,colorbar,
    colorbar style={colormap access=const}
]
    \addplot3[surf,shader=interp,colormap access=const] {x*y};
\end{axis}
\end{tikzpicture}
\end{codeexample}
            %
            注意，最后一个示例的外观与 |contour filled| 生成的相同。
            区别在于 |contour filled| 接受输入颜色映射，并根据所选
            等高线层级\emph{重新采样}它。上面的示例只是一个使用特殊
            颜色映射的``普通''曲面图。
        \end{pgfplotskeylist}
    \item
        \declareandlabel{const color of colormap}|(|\meta{position}|)=(|\meta{value}| of |\meta{colormap name}|)|%
            \index{const color of colormap!In colormap specification}%

        该键与上面提到的 |color of colormap| 几乎相同，但它使用与
        |colormap access=piecewise constant| 相同的功能，同时从\emph{源}
        颜色映射确定颜色（包括 |colormap access/extra interval width|）。
        注意，生成的颜色映射仍可与 |colormap access| 的任何值配合使用，
        包括 |colormap access=const| 和 |colormap access=map|。

        注意，|const color of colormap| 也可用作绘图操作的键：

        \begin{pgfplotskey}{const color of colormap=\meta{value} of \meta{colormap name}}
           使用特殊颜色名 `|.|' 选择 \meta{colormap name} 的指定颜色，
           并赋值 |color=.|：
            %
\begin{codeexample}[]
\pgfplotsset{
  colormap={example}{color=(red) color=(blue) color=(yellow)}}
\tikz \fill[const color of colormap=50 of example] (0,0) rectangle (1,1);
\tikz \fill[const color of colormap=750 of example] (0,0) rectangle (1,1);
\end{codeexample}
        \end{pgfplotskey}

        |samples of colormap| 的所有特殊说明（如花括号、选项列表支持、
        位置）也适用于此。
    \item
        \declareandlabel{const colors of colormap}|(|\meta{position}|)=(|\meta{list}| of |\meta{colormap name}|)|%
            \index{const colors of colormap!In colormap specification}%

        一个便捷键，等价于 |const color of colormap(|\meta{position}|)| 序列，
        \meta{list} 中每个元素一个。\meta{list} 使用 |\foreach| 求值。

        注意，|const colors of colormap| 在 |cycle list| 定义期间也可用。

        |samples of colormap| 的所有特殊说明（如花括号、选项列表支持、
        位置）也适用于此。
\end{enumerate}

\begin{pgfplotskey}{colormap access/extra interval width=\marg{fraction} (initially h)}
    \paragraph{注意：}

    此键应是实现的技术部分。你可能想跳过其文档，因为它通常开箱即用。
    你只需注意 |sample for|。

    此键\emph{仅}适用于 |colormap access=piecewise constant|：它确保
    |colormap| 中的每种颜色获得自己的\emph{区间}。这确保每种颜色在输出中
    实际可见。因此，每种提供的颜色类似一个\emph{区间}。这与
    |colormap access=map| 不同，其中每种提供的颜色类似一个\emph{区间边界}。

    通常，\PGFPlots{} 仅在自动计算位置时激活此功能。因此，以下示例隐式
    使用 |extra interval width=h|（默认）：
    %
\begin{codeexample}[]
\pgfplotscolorbardrawstandalone[
    colormap={example}{
        color=(blue)
        color=(red)
        color=(black)
    },
    colorbar horizontal,
    colorbar style={xtick=data},
    colormap access=const,
]
\end{codeexample}

    一旦手动提供位置，\PGFPlots{} 默认使用
    |extra interval width=0| 以尊重输入设置：
    %
\begin{codeexample}[]
\pgfplotscolorbardrawstandalone[
    colormap={example}{
        color(0)=(blue)
        color(500)=(red)
        color(1000)=(black)
    },
    colorbar horizontal,
    colorbar style={xtick=data},
    colormap access=const,
]
\end{codeexample}
    %
    \noindent 注意，位置构成了 |colormap| 的输入节点，其结果是最后一种
    颜色\emph{未被使用}；它仅用作区间边界。在此上下文中，``手动提供''
    指圆括号中的位置或 |target pos| 的非空值。

    要覆盖输入位置并获得额外区间，你必须设置选项：
    %
\begin{codeexample}[]
\pgfplotscolorbardrawstandalone[
    colormap={example}{
        color(0)=(blue)
        color(500)=(red)
        color(1000)=(black)
    },
    colorbar horizontal,
    colorbar style={xtick=data},
    colormap access/extra interval width=h,
    colormap access=const,
]
\end{codeexample}
    %
    注意 |colormap access=const| 的额外区间如何修改颜色映射定义的位置：
    它们都向左移动，所选输入位置相应缩放。下面的示例中这一点更明显。
    首先，我们用默认设置生成一个颜色映射，它禁用额外区间，但省略最后
    （最亮）的颜色：
    %
\begin{codeexample}[]
\pgfplotscolorbardrawstandalone[
    colormap={nonuniform}{
        of colormap={
            viridis,
            target pos={0,200,300,500,700,1000}
        },
    },
    colorbar horizontal,
    colormap access=const,
    colorbar style={xtick=data,font=\tiny,
        /pgf/number format/precision=0},
    colormap access=const,
]
\end{codeexample}

    接下来，我们显式启用额外区间，看到输入位置向左缩放。
    注意它们保持相对距离，但颜色映射的最后一种颜色终于可见：
    %
\begin{codeexample}[]
\pgfplotscolorbardrawstandalone[
    colormap={nonuniform}{
        of colormap={
            viridis,
            target pos={0,200,300,500,700,1000}
        },
    },
    colorbar horizontal,
    colormap access=const,
    colormap access/extra interval width=h,
    colorbar style={xtick=data,font=\tiny,
        /pgf/number format/precision=0},
    colormap access=const,
]
\end{codeexample}

    最后一个示例依赖 |of colormap| 和\emph{采样}过程。在此上下文中，
    \PGFPlots{} 提供 |sample for=const|，产生预期外观：
    %
\begin{codeexample}[]
\pgfplotscolorbardrawstandalone[
    colormap={nonuniform}{
        of colormap={
            viridis,
            target pos={0,200,300,500,700,1000},
            sample for=const,
        },
    },
    colorbar horizontal,
    colormap access=const,
    colorbar style={xtick=data,font=\tiny,
        /pgf/number format/precision=0},
    colormap access=const,
]
\end{codeexample}
    %
    请参阅 |sample for| 的文档。

    该区间的默认宽度是颜色映射的网格宽度（值 \declaretext{h}），
    即它始终与其他区间一样宽。如果颜色映射有非均匀距离，
    则使用遇到的最小网格宽度作为额外区间。如果我们省略该键并在右端
    手动添加一些人工颜色，就可以看到这一点——我们只需确保最右端区间
    长度正确。在上面的示例中，最小网格宽度为 $100$，因此我们可以通过
    以下方式生成等效结果
    %
    \index{Color maps!Non uniform}
\begin{codeexample}[]
\pgfplotscolorbardrawstandalone[
    colormap={nonuniform}{
        of colormap={
            viridis,
            target pos={0,200,300,500,700,1000}
        },
        color(1100)=(red)
    },
    colorbar horizontal,
    colormap access=const,
    colorbar style={xtick=data,font=\tiny,
        /pgf/number format/precision=0},
    colormap access=const,
]
\end{codeexample}
    %
    \noindent 如前所述，|colormap access=const| 忽略最右端颜色（``|red|''）。
    注意，此方法几乎与 |sample for=const| 的内部实现相同。

    |colormap access/extra interval width=|\meta{fraction} 的值也可用于
    自定义人工区间的宽度：它是总宽度的分数，接受值为
    $0\le$\meta{fraction}$\le 0.9$，其中 $0$ 禁用额外区间，
    $0.9$ 对应结果宽度的 $90\%$。\meta{fraction} 的任何非 $0$ 值都会
    创建一个额外区间，其宽度为整个颜色映射的 \meta{fraction} 百分比。
    特殊魔术值 |colormap access/extra interval width=|\declaretext{h}
    将使用颜色映射的网格宽度。如果颜色映射有非均匀距离，它将使用遇到的
    最小网格宽度。这是默认值。

    注意，\meta{fraction} 是 |colormap| 的属性。该键仅为\emph{当前}颜色映射
    定义它。定义新的 |colormap| 时，新 |colormap| 使用默认宽度
    （但保留旧 |colormap| 的配置值）。


    \paragraph{注意：}

    在撰写本文时，均匀颜色映射仅支持默认区间宽度 `|h|' 和
    \meta{fraction}$=0$，进一步自定义仅适用于非均匀颜色映射。
    如果你需要绕过此限制，应提交功能请求，并移动某个颜色位置，
    直到颜色映射变为非均匀颜色映射。
\end{pgfplotskey}

}%


\subsubsection{将颜色映射条目选为普通颜色}

\begin{pgfplotskeylist}{%
    color of colormap=\meta{value},
    color of colormap=\meta{value} of \meta{colormap name}%
}
    [另见第~\ref{sec:pgfplots:colormaps:based:on:others} 节
    （第~\pageref{sec:pgfplots:colormaps:based:on:others} 页），
    了解如何在 |colormap| 定义中使用它]

    将 \tikzname{} |color| 定义为 \meta{colormap name} 的 \meta{value}。
    如果省略 \meta{colormap name}，则求值 |colormap name| 的值
    （即使用当前颜色映射）。

    参数 \meta{value} 应为范围 $[0,1000]$ 内的数字，其中 $0$ 表示颜色映射
    的下端，$1000$ 表示上端。

    \pgfkeysgetvalue{/pgfplots/colormap name}\colormapname
    当前颜色映射（\colormapname）：\pgfplotsshowcolormap{\colormapname}

\begin{codeexample}[]
\tikz\fill [color of colormap={800}, thick,
    draw=.!60!black]
        (0,0) -- (1,1) -- (2,1) circle (10pt);
\end{codeexample}

    该键计算请求的颜色并调用 |color=.|。请记住，魔术颜色名 `|.|'
    始终反映``当前颜色''，即 |color=|\meta{some color} 的结果。
    还要记住，|color| 是 \tikzname{} 命令，它仅定义颜色，
    你还必须提供 `|draw|' 或 `|fill|' 之一才能生效。由于 `|.|'
    是普通颜色，我们可以写 |draw=.!60!black| 将其与另一种颜色组合。

    viridis: \pgfplotsshowcolormap{viridis}

\begin{codeexample}[]
\pgfplotsset{colormap name=viridis}
\tikz\fill [color of colormap={1000},thick,
    draw=.!60!black]
        (0,0) -- (1,1) -- (2,1) circle (10pt);
\end{codeexample}

\begin{codeexample}[]
\tikz\fill [color of colormap={300 of viridis},thick,
    draw=.!60!black]
        (0,0) -- (1,1) -- (2,1) circle (10pt);
\end{codeexample}

    参数 \meta{colormap name} 既可以是 |colormap name| 的有效参数，
    \emph{也}可以是样式名，如 |colormap/cool|：

    cool: \pgfplotsshowcolormap{cool}

\begin{codeexample}[]
\tikz\fill [color of colormap={300 of colormap/cool},
    thick,draw=.!60!black]
        (0,0) -- (1,1) -- (2,1) circle (10pt);
\end{codeexample}

    后一种语法允许惰性求值颜色映射。但是，如果你多次引用同一颜色映射，先写
    |\pgfplotsset{colormap/cool}|，再多次引用
    |color of colormap={... of cool}| 是有意义的，以避免不必要的
    惰性求值。

    可以编写大量不带显式 \meta{colormap name} 的调用，即大量
    |color of colormap=|\meta{value} 形式的调用。它们都将使用当时
    激活的 |colormap name|。

    注意，实际上有两个键前缀的键：|/pgfplots/color of colormap|
    和别名 |/tikz/color of colormap|。这允许将这些键同时用于普通
    \tikzname{} 图形和 \PGFPlots{}。

    另见 |colormap access=map|。
\end{pgfplotskeylist}

\begin{pgfplotskeylist}{%
    index of colormap=\meta{index},
    index of colormap=\meta{index} of \meta{colormap name}%
}
    [另见第~\ref{sec:pgfplots:colormaps:based:on:others} 节
    （第~\pageref{sec:pgfplots:colormaps:based:on:others} 页），
    了解如何在 |colormap| 定义中使用它]

    |color of colormap| 的一个变体，它按索引访问 \meta{colormap name}。
    因此，参数 \meta{index} 是范围 $0,\dotsc,N-1$ 内的整数，
    其中 $N$ 是定义 \meta{colormap name} 的颜色数量。
    超出此范围的 \meta{index} 会自动裁剪到上界。

\begin{codeexample}[]
    \pgfplotsset{colormap/jet}
\foreach \i in {
    0,...,\pgfplotscolormaplastindexof{jet}
}{
    \tikz\fill [
        index of colormap={\i of jet},
        thick,
        draw=.!60!black,
    ] (0,0) rectangle (10pt,6pt);
}
\end{codeexample}

    \begin{command}{\pgfplotscolormapsizeof\marg{colormap name}}
        展开为构成 \meta{colormap name} 的颜色数量。

        如果参数 \meta{colormap name} 是未知颜色映射，则展开为 $0$。
    \end{command}

    \begin{command}{\pgfplotscolormaplastindexof\marg{colormap name}}
        展开为 \meta{colormap name} 的最后一个索引，即它是访问
        $N-1$ 的便捷方法。

        如果参数 \meta{colormap name} 是未知颜色映射，则展开为 $-1$。
    \end{command}

    另见 |colormap access=direct|。
\end{pgfplotskeylist}

\begin{pgfplotskeylist}{%
    const color of colormap=\meta{value},
    const color of colormap=\meta{value} of \meta{colormap name}%
}
    [另见第~\ref{sec:pgfplots:colormaps:based:on:others} 节
    （第~\pageref{sec:pgfplots:colormaps:based:on:others} 页），
    了解如何在 |colormap| 定义中使用它]

    将 \tikzname{} |color| 定义为 \meta{colormap name} 的 \meta{value}。
    如果省略 \meta{colormap name}，则求值 |colormap name| 的值
    （即使用当前颜色映射）。

    该键与 |color of colormap| 几乎相同，只是它使用
    |colormap access=piecewise constant| 来确定插值后的值。
\end{pgfplotskeylist}


\subsection{循环列表——控制线样式的选项}
\label{sec:cycle:list}

\begin{pgfplotskeylist}{cycle list=\marg{list},cycle list name=\marg{name} (initially color)}
    允许指定一系列绘图说明，用于每个没有显式绘图说明的
    \hbox{|\addplot|} 命令。因此，如果你写
    |\addplot+|\oarg{keys}| ...;|，或者像
    |\addplot|\oarg{explicit plot specification}| ...;| 那样\emph{不}使用
    方括号，就会使用当前激活的 |cycle list|。

    将选择索引为~$i$ 的列表元素，其中~$i$ 是当前 |\addplot| 命令的索引
    （另见 |cycle list shift| 键，它允许改用 $i+n$）。此索引也包括
    不使用 |cycle list| 的绘图命令。

    有多种方法可以更改当前激活的 |cycle list|：


    \subsubsection{预定义列表}

    使用预定义列表之一，\footnote{在早期版本中，这些列表称为
    \texttt{\textbackslash coloredplotspeclist} 和
    \texttt{\textbackslash blackwhiteplotspeclist}，显得过长，
    因此已重命名。不过旧名称仍被接受。}
    %
    \begin{itemize}
        \item \declareandlabel{color}（从上到下）
            %
\begin{codeexample}[]
\begin{tikzpicture}
\begin{axis}[
    stack plots=y,stack dir=minus,
    cycle list name=color,
]
    \addplot coordinates {(0,1) (0.5,1) (1,1)};
    \addplot coordinates {(0,1) (0.5,1) (1,1)};
    \addplot coordinates {(0,1) (0.5,1) (1,1)};
    \addplot coordinates {(0,1) (0.5,1) (1,1)};
    \addplot coordinates {(0,1) (0.5,1) (1,1)};
    \addplot coordinates {(0,1) (0.5,1) (1,1)};
    \addplot coordinates {(0,1) (0.5,1) (1,1)};
    \addplot coordinates {(0,1) (0.5,1) (1,1)};
    \addplot coordinates {(0,1) (0.5,1) (1,1)};
    \addplot coordinates {(0,1) (0.5,1) (1,1)};
    \addplot coordinates {(0,1) (0.5,1) (1,1)};
    \addplot coordinates {(0,1) (0.5,1) (1,1)};
    \addplot coordinates {(0,1) (0.5,1) (1,1)};
\end{axis}
\end{tikzpicture}
\end{codeexample}
            %
        \item \declareandlabel{exotic} (from top to bottom)
            %
\begin{codeexample}[]
\begin{tikzpicture}
\begin{axis}[
    stack plots=y,stack dir=minus,
    cycle list name=exotic,
]
    \addplot coordinates {(0,1) (0.5,1) (1,1)};
    \addplot coordinates {(0,1) (0.5,1) (1,1)};
    \addplot coordinates {(0,1) (0.5,1) (1,1)};
    \addplot coordinates {(0,1) (0.5,1) (1,1)};
    \addplot coordinates {(0,1) (0.5,1) (1,1)};
    \addplot coordinates {(0,1) (0.5,1) (1,1)};
    \addplot coordinates {(0,1) (0.5,1) (1,1)};
    \addplot coordinates {(0,1) (0.5,1) (1,1)};
    \addplot coordinates {(0,1) (0.5,1) (1,1)};
    \addplot coordinates {(0,1) (0.5,1) (1,1)};
    \addplot coordinates {(0,1) (0.5,1) (1,1)};
    \addplot coordinates {(0,1) (0.5,1) (1,1)};
    \addplot coordinates {(0,1) (0.5,1) (1,1)};
\end{axis}
\end{tikzpicture}
\end{codeexample}

        \item \declareandlabel{black white} (from top to bottom)
            %
\begin{codeexample}[]
\begin{tikzpicture}
\begin{axis}[
    stack plots=y,stack dir=minus,
    cycle list name=black white,
]
    \addplot coordinates {(0,1) (0.5,1) (1,1)};
    \addplot coordinates {(0,1) (0.5,1) (1,1)};
    \addplot coordinates {(0,1) (0.5,1) (1,1)};
    \addplot coordinates {(0,1) (0.5,1) (1,1)};
    \addplot coordinates {(0,1) (0.5,1) (1,1)};
    \addplot coordinates {(0,1) (0.5,1) (1,1)};
    \addplot coordinates {(0,1) (0.5,1) (1,1)};
    \addplot coordinates {(0,1) (0.5,1) (1,1)};
    \addplot coordinates {(0,1) (0.5,1) (1,1)};
    \addplot coordinates {(0,1) (0.5,1) (1,1)};
    \addplot coordinates {(0,1) (0.5,1) (1,1)};
    \addplot coordinates {(0,1) (0.5,1) (1,1)};
    \addplot coordinates {(0,1) (0.5,1) (1,1)};
\end{axis}
\end{tikzpicture}
\end{codeexample}
            %
        \item \declareandlabel{mark list} (from top to bottom)
            %
\begin{codeexample}[]
\begin{tikzpicture}
\begin{axis}[
    stack plots=y,stack dir=minus,
    cycle list name=mark list,
]
  \addplot+ [blue] coordinates {(0,1) (0.5,1) (1,1)};
  \addplot+ [blue] coordinates {(0,1) (0.5,1) (1,1)};
  \addplot+ [blue] coordinates {(0,1) (0.5,1) (1,1)};
  \addplot+ [blue] coordinates {(0,1) (0.5,1) (1,1)};
  \addplot+ [blue] coordinates {(0,1) (0.5,1) (1,1)};
  \addplot+ [blue] coordinates {(0,1) (0.5,1) (1,1)};
  \addplot+ [blue] coordinates {(0,1) (0.5,1) (1,1)};
  \addplot+ [blue] coordinates {(0,1) (0.5,1) (1,1)};
  \addplot+ [blue] coordinates {(0,1) (0.5,1) (1,1)};
  \addplot+ [blue] coordinates {(0,1) (0.5,1) (1,1)};
  \addplot+ [blue] coordinates {(0,1) (0.5,1) (1,1)};
  \addplot+ [blue] coordinates {(0,1) (0.5,1) (1,1)};
  \addplot+ [blue] coordinates {(0,1) (0.5,1) (1,1)};
\end{axis}
\end{tikzpicture}
\end{codeexample}

            |mark list| 始终使用当前颜色，但不定义颜色
            （\verbpdfref{\addplot+} 语句显式将当前颜色设为 |blue|）。

            |mark list| 在与 |cycle multi list| 配合时特别有用，
            它允许将其与其他列表（例如 |linestyles| 或颜色列表）组合。
        \item \declareandlabel{mark list*} 一个仅含标记的列表。
            与 |mark list| 不同，所有这些标记都是填充的。
            它们定义为（从上到下）
            %
\begin{codeexample}[]
\begin{tikzpicture}
\begin{axis}[
    stack plots=y,stack dir=minus,
    cycle list name=mark list*,
]
  \addplot+ [blue] coordinates {(0,1) (0.5,1) (1,1)};
  \addplot+ [blue] coordinates {(0,1) (0.5,1) (1,1)};
  \addplot+ [blue] coordinates {(0,1) (0.5,1) (1,1)};
  \addplot+ [blue] coordinates {(0,1) (0.5,1) (1,1)};
  \addplot+ [blue] coordinates {(0,1) (0.5,1) (1,1)};
  \addplot+ [blue] coordinates {(0,1) (0.5,1) (1,1)};
  \addplot+ [blue] coordinates {(0,1) (0.5,1) (1,1)};
  \addplot+ [blue] coordinates {(0,1) (0.5,1) (1,1)};
  \addplot+ [blue] coordinates {(0,1) (0.5,1) (1,1)};
  \addplot+ [blue] coordinates {(0,1) (0.5,1) (1,1)};
  \addplot+ [blue] coordinates {(0,1) (0.5,1) (1,1)};
  \addplot+ [blue] coordinates {(0,1) (0.5,1) (1,1)};
  \addplot+ [blue] coordinates {(0,1) (0.5,1) (1,1)};
\end{axis}
\end{tikzpicture}
\end{codeexample}
            %
            与 |mark list| 类似，|mark list*| 始终使用当前颜色，
            但不定义颜色（\verbpdfref{\addplot+} 见上文）。
        \item \declareandlabel{color list}（从上到下）
            %
\begin{codeexample}[]
\begin{tikzpicture}
\begin{axis}[
    stack plots=y,stack dir=minus,
    cycle list name=color list,
]
    \addplot coordinates {(0,1) (0.5,1) (1,1)};
    \addplot coordinates {(0,1) (0.5,1) (1,1)};
    \addplot coordinates {(0,1) (0.5,1) (1,1)};
    \addplot coordinates {(0,1) (0.5,1) (1,1)};
    \addplot coordinates {(0,1) (0.5,1) (1,1)};
    \addplot coordinates {(0,1) (0.5,1) (1,1)};
    \addplot coordinates {(0,1) (0.5,1) (1,1)};
    \addplot coordinates {(0,1) (0.5,1) (1,1)};
    \addplot coordinates {(0,1) (0.5,1) (1,1)};
    \addplot coordinates {(0,1) (0.5,1) (1,1)};
    \addplot coordinates {(0,1) (0.5,1) (1,1)};
    \addplot coordinates {(0,1) (0.5,1) (1,1)};
    \addplot coordinates {(0,1) (0.5,1) (1,1)};
\end{axis}
\end{tikzpicture}
\end{codeexample}

            选项 |cycle list name=color| 也使用标记，而
            |color list| \emph{仅}使用颜色。

        \item \declareandlabel{linestyles}（从上到下）
            %
\begin{codeexample}[]
\begin{tikzpicture}
\begin{axis}[
    stack plots=y,stack dir=minus,
    cycle list name=linestyles,
]
    \addplot coordinates {(0,1) (0.5,1) (1,1)};
    \addplot coordinates {(0,1) (0.5,1) (1,1)};
    \addplot coordinates {(0,1) (0.5,1) (1,1)};
    \addplot coordinates {(0,1) (0.5,1) (1,1)};
    \addplot coordinates {(0,1) (0.5,1) (1,1)};
    \addplot coordinates {(0,1) (0.5,1) (1,1)};
    \addplot coordinates {(0,1) (0.5,1) (1,1)};
    \addplot coordinates {(0,1) (0.5,1) (1,1)};
    \addplot coordinates {(0,1) (0.5,1) (1,1)};
    \addplot coordinates {(0,1) (0.5,1) (1,1)};
    \addplot coordinates {(0,1) (0.5,1) (1,1)};
    \addplot coordinates {(0,1) (0.5,1) (1,1)};
    \addplot coordinates {(0,1) (0.5,1) (1,1)};
\end{axis}
\end{tikzpicture}
\end{codeexample}
            %
        \item \declareandlabel{linestyles*} 比 |linestyles| 包含更多
            点划线样式（从上到下）
            %
\begin{codeexample}[]
\begin{tikzpicture}
\begin{axis}[
    stack plots=y,stack dir=minus,
    cycle list name=linestyles*,
]
    \addplot coordinates {(0,1) (0.5,1) (1,1)};
    \addplot coordinates {(0,1) (0.5,1) (1,1)};
    \addplot coordinates {(0,1) (0.5,1) (1,1)};
    \addplot coordinates {(0,1) (0.5,1) (1,1)};
    \addplot coordinates {(0,1) (0.5,1) (1,1)};
    \addplot coordinates {(0,1) (0.5,1) (1,1)};
    \addplot coordinates {(0,1) (0.5,1) (1,1)};
    \addplot coordinates {(0,1) (0.5,1) (1,1)};
    \addplot coordinates {(0,1) (0.5,1) (1,1)};
    \addplot coordinates {(0,1) (0.5,1) (1,1)};
    \addplot coordinates {(0,1) (0.5,1) (1,1)};
    \addplot coordinates {(0,1) (0.5,1) (1,1)};
    \addplot coordinates {(0,1) (0.5,1) (1,1)};
\end{axis}
\end{tikzpicture}
\end{codeexample}
        \item \declareandlabel{auto} |cycle list name=auto| 始终表示
            最近由 |cycle list| 或 |cycle list name| 激活的循环列表。
    \end{itemize}

    所有预定义循环列表的定义如下（语法说明见本段末尾）。
        %
\begin{codeexample}[code only]
\pgfplotscreateplotcyclelist{color}{
    blue,every mark/.append style={fill=blue!80!black},mark=*\\
    red,every mark/.append style={fill=red!80!black},mark=square*\\
    brown!60!black,every mark/.append style={fill=brown!80!black},mark=otimes*\\
    black,mark=star\\
    blue,every mark/.append style={fill=blue!80!black},mark=diamond*\\
    red,densely dashed,every mark/.append style={solid,fill=red!80!black},mark=*\\
    brown!60!black,densely dashed,every mark/.append style={
        solid,fill=brown!80!black},mark=square*\\
    black,densely dashed,every mark/.append style={solid,fill=gray},mark=otimes*\\
    blue,densely dashed,mark=star,every mark/.append style=solid\\
    red,densely dashed,every mark/.append style={solid,fill=red!80!black},mark=diamond*\\
}
\end{codeexample}

\begin{codeexample}[code only]
\pgfplotscreateplotcyclelist{black white}{
    every mark/.append style={fill=gray},mark=*\\
    every mark/.append style={fill=gray},mark=square*\\
    every mark/.append style={fill=gray},mark=otimes*\\
    mark=star\\
    every mark/.append style={fill=gray},mark=diamond*\\
    densely dashed,every mark/.append style={solid,fill=gray},mark=*\\
    densely dashed,every mark/.append style={solid,fill=gray},mark=square*\\
    densely dashed,every mark/.append style={solid,fill=gray},mark=otimes*\\
    densely dashed,every mark/.append style={solid},mark=star\\
    densely dashed,every mark/.append style={solid,fill=gray},mark=diamond*\\
}
\end{codeexample}

\begin{codeexample}[code only]
\pgfplotscreateplotcyclelist{exotic}{
    teal,every mark/.append style={fill=teal!80!black},mark=*\\
    orange,every mark/.append style={fill=orange!80!black},mark=square*\\
    cyan!60!black,every mark/.append style={fill=cyan!80!black},mark=otimes*\\
    red!70!white,mark=star\\
    lime!80!black,every mark/.append style={fill=lime},mark=diamond*\\
    red,densely dashed,every mark/.append style={solid,fill=red!80!black},mark=*\\
    yellow!60!black,densely dashed,
        every mark/.append style={solid,fill=yellow!80!black},mark=square*\\
    black,every mark/.append style={solid,fill=gray},mark=otimes*\\
    blue,densely dashed,mark=star,every mark/.append style=solid\\
    red,densely dashed,every mark/.append style={solid,fill=red!80!black},mark=diamond*\\
}
\end{codeexample}

\begin{codeexample}[code only]
% note that "." is the currently defined Tikz color.
\pgfplotscreateplotcyclelist{mark list}{
    every mark/.append style={solid,fill=\pgfplotsmarklistfill},mark=*\\
    every mark/.append style={solid,fill=\pgfplotsmarklistfill},mark=square*\\
    every mark/.append style={solid,fill=\pgfplotsmarklistfill},mark=triangle*\\
    every mark/.append style={solid},mark=star\\
    every mark/.append style={solid,fill=\pgfplotsmarklistfill},mark=diamond*\\
    every mark/.append style={solid,fill=\pgfplotsmarklistfill!40},mark=otimes*\\
    every mark/.append style={solid},mark=|\\
    every mark/.append style={solid,fill=\pgfplotsmarklistfill},mark=pentagon*\\
    every mark/.append style={solid},mark=text,text mark=p\\
    every mark/.append style={solid},mark=text,text mark=a\\
}
\end{codeexample}
    %
    \noindent 在此上下文中，通用填充颜色表达式可使用
    |mark list fill| 自定义：
    %
    \begin{pgfplotskey}{mark list fill=\marg{color} (initially .!80!black)}
        允许自定义 |mark list| 和 |mark list*| 的填充颜色。

        例如，如果颜色为 |black|，替代选项
        |mark list fill=.!50!white| 会产生好得多的效果。
    \end{pgfplotskey}

    \begin{command}{\pgfplotsmarklistfill}
        展开为 |\pgfkeysvalueof{||/pgfplots/mark list fill||}|。
    \end{command}
    %
\begin{codeexample}[code only]
% note that "." is the currently defined Tikz color.
\pgfplotscreateplotcyclelist{mark list*}{
    every mark/.append style={solid,fill=\pgfplotsmarklistfill},mark=*\\
    every mark/.append style={solid,fill=\pgfplotsmarklistfill},mark=square*\\
    every mark/.append style={solid,fill=\pgfplotsmarklistfill},mark=triangle*\\
    every mark/.append style={solid,fill=\pgfplotsmarklistfill},mark=halfsquare*\\
    every mark/.append style={solid,fill=\pgfplotsmarklistfill},mark=pentagon*\\
    every mark/.append style={solid,fill=\pgfplotsmarklistfill},mark=halfcircle*\\
    every mark/.append style={solid,fill=\pgfplotsmarklistfill,rotate=180},mark=halfdiamond*\\
    every mark/.append style={solid,fill=\pgfplotsmarklistfill!40},mark=otimes*\\
    every mark/.append style={solid,fill=\pgfplotsmarklistfill},mark=diamond*\\
    every mark/.append style={solid,fill=\pgfplotsmarklistfill},mark=halfsquare right*\\
    every mark/.append style={solid,fill=\pgfplotsmarklistfill},mark=halfsquare left*\\
}
\end{codeexample}

\begin{codeexample}[code only]
\pgfplotscreateplotcyclelist{color list}{
    red,blue,black,yellow,brown,teal,orange,violet,cyan,green!70!black,magenta,gray
}
\end{codeexample}

\begin{codeexample}[code only]
\pgfplotscreateplotcyclelist{linestyles}{solid,dashed,dotted}
\pgfplotscreateplotcyclelist{linestyles*}{solid,dashed,dotted,dashdotted,dashdotdotted}
\end{codeexample}


    \subsubsection{定义自己的循环列表}

    循环列表的第二种选择是直接将每个条目作为 |cycle list| 的参数提供，
    %
\begin{codeexample}[]
\begin{tikzpicture}
\begin{loglogaxis}[
    cycle list={
        {blue,mark=*},
        {red,mark=square},
        {dashed,mark=o},
        {loosely dotted,mark=+},
        {brown!60!black,mark=otimes*,
            mark options={fill=brown!40},
        }% <-- don't add a comma here
    },
]
    \plotcoords
    \legend{$d=2$,$d=3$,$d=4$,$d=5$,$d=6$}
\end{loglogaxis}
\end{tikzpicture}
\end{codeexample}
    %
    （此示例列表需要 |\usetikzlibrary{plotmarks}|）。

    输入格式在下文更详细地描述。


    \subsubsection{定义并命名自己的循环列表}

    循环列表的最后一种方法是在导言区组合定义\emph{命名}循环列表，
    并使用 `|cycle list name|' 使用它们：
    %
    \begin{command}{\pgfplotscreateplotcyclelist\allowbreak\marg{name}\allowbreak\marg{list}}%
    \end{command}
    %
\begin{codeexample}[code only]
\pgfplotscreateplotcyclelist{mylist}{
    {blue,mark=*},
    {red,mark=square},
    {dashed,mark=o},
    {loosely dotted,mark=+},
    {brown!60!black,mark options={fill=brown!40},mark=otimes*}% <-- don't add a comma here
}
...
\begin{axis}[cycle list name=mylist]
    ...
\end{axis}
\end{codeexample}


    \subsubsection{定义循环列表：输入格式}

    |cycle list| 由形如 |cycle list=|\marg{list} 的键值对或上述等效宏
    |\pgfplotscreateplotcyclelist|\allowbreak\marg{name}\allowbreak\marg{list} 定义。

    在此上下文中，参数 \meta{list} 通常是以逗号分隔的样式键列表，
    如颜色、线型、标记类型和标记样式。这种``逗号列表的逗号列表''结构
    需要用花括号封装内部列表：
    %
\begin{codeexample}[code only]
\pgfplotscreateplotcyclelist{mylist}{
    {blue,mark=*},
    {red,mark=square},
    {dashed,mark=o},
    {loosely dotted,mark=+},
    {brown!60!black,mark options={fill=brown!40},mark=otimes*}% <-- don't add a comma here
}
\end{codeexample}
    %
    或者，可以用 `|\\|' 终止内部列表（即单个绘图的列表）：
    %
\begin{codeexample}[code only]
\begin{axis}[
    cycle list={
        blue,mark=*\\
        red,mark=square\\
        dashed,mark=o\\
        loosely dotted,mark=+\\
        brown!60!black,mark options={fill=brown!40},mark=otimes*\\
    },
]
...
\end{axis}
\end{codeexample}
    %
    此时，\emph{最后一个}条目也需要终止符 `|\\|'，但可以省略单个条目
    周围的花括号。


    \subsubsection{操作循环列表与绘图的关联}

    \begin{pgfplotskey}{cycle list shift=\marg{integer} (initially empty)}
        允许\emph{偏移} |cycle list| 中的索引。如果 \meta{integer} 为 $n$，
        则取列表元素 $i+n$ 而非第 $i$ 个。记住 $i$ 是当前 |\addplot|
        命令的索引（从~$0$ 开始）。

        由于在调用 |\addplot|（或 |\addplot+|）时\emph{立即}查询
        |cycle list|，你可以为选定的绘图调整 |cycle list shift|：
        %
\begin{codeexample}[code only]
    \pgfplotsset{cycle list shift=3}
\addplot ...

    \pgfplotsset{cycle list shift=-1}
\addplot ...
\end{codeexample}

        \paragraph{特殊情况：}

        如果结果为负，即 $i+n<0$，则取列表索引 $-(i+n)$。例如，
        |cycle list shift=-10| 且 $i<10$ 将得到列表索引 $10-i$。
        注意，你可以使用 |reverse legend| 反转图例，因此此功能
        可能永远不需要。
    \end{pgfplotskey}


    \subsubsection{基于颜色映射定义循环列表}

        除了从头定义 |cycle list| 外，\PGFPlots{} 还支持
        |cycle list=|\meta{list} 的专用输入定义，允许从现有
        |colormap| 获取值。此时，\meta{list} 包含方括号括起的键：
    \index{Colormap!Convert into Cycle list}

{%
\pgfplotsset{
    % ATTENTION: this is documented in this manual!
    cycle from colormap manual style/.style={
        x=3cm,y=10pt,ytick=\empty,
        colorbar style={x=,y=,ytick=\empty},
        point meta min=0,point meta max=1,
        stack plots=y,
        y dir=reverse,colorbar style={y dir=reverse},
        every axis plot/.style={line width=2pt},
        legend entries={0,...,20},
        legend pos=outer north east,
    },
}

    第一种语法 \declareandlabel{of colormap} 允许将 |colormap| 的颜色转换为
    |cycle list|。它可通过
    |cycle list={[of colormap]}| 无参数指定，以取
    最近分配的 |colormap name| 的值（即当前 |colormap|）。它也可指定为
    |cycle list={[of colormap=name]}|，此时使用指定的
    |colormap name=name|。
    在两种情况下，定义仅将颜色映射中找到的颜色转换为循环列表，即不涉及
    插值。将此应用于默认 |colormap name=hot|（有~$4$ 种颜色）得到以下示例：
    %
\begin{codeexample}[]
\begin{tikzpicture}
\begin{axis}[
    colormap name=hot,colorbar,
    cycle list={[of colormap]},
    %
    cycle from colormap manual style,
]
    \addplot coordinates {(0,1) (0.5,1) (1,1)};
    \addplot coordinates {(0,1) (0.5,1) (1,1)};
    \addplot coordinates {(0,1) (0.5,1) (1,1)};
    \addplot coordinates {(0,1) (0.5,1) (1,1)};
    \addplot coordinates {(0,1) (0.5,1) (1,1)};
    \addplot coordinates {(0,1) (0.5,1) (1,1)};
    \addplot coordinates {(0,1) (0.5,1) (1,1)};
    \addplot coordinates {(0,1) (0.5,1) (1,1)};
    \addplot coordinates {(0,1) (0.5,1) (1,1)};
    \addplot coordinates {(0,1) (0.5,1) (1,1)};
    \addplot coordinates {(0,1) (0.5,1) (1,1)};
\end{axis}
\end{tikzpicture}
\end{codeexample}
    %
    注意，由于 |hot| 有 $4$ 种颜色，|cycle list| 也包含 $4$ 个条目，
    每 $4$ 个绘图重复一次。

    第二种方式类似于 |samples of colormap|：它接受
    \declaretext{samples of colormap}|=|\marg{number}\allowbreak| of |\meta{colormap name}。
    它选择所选颜色映射的 \meta{number} 个样本。

    第三种方式类似于 |color of colormap|：它接受\linebreak
    |colors of colormap =|\marg{list} 或\linebreak
    |colors of colormap=|\marg{list}\allowbreak| of |\meta{colormap name}。
    该选项\emph{插值}颜色，并接受范围 $[0,1000]$ 内的值 \meta{list}，
    其中 $0$ 是 |colormap| 的最低元素，$1000$ 是其最高元素：
    %
\begin{codeexample}[]
\begin{tikzpicture}
\begin{axis}[
    colormap name=hot,colorbar,
    cycle list={
        [colors of colormap={0,100,...,1000}]},
    %
    cycle from colormap manual style,
]
    \addplot coordinates {(0,1) (0.5,1) (1,1)};
    \addplot coordinates {(0,1) (0.5,1) (1,1)};
    \addplot coordinates {(0,1) (0.5,1) (1,1)};
    \addplot coordinates {(0,1) (0.5,1) (1,1)};
    \addplot coordinates {(0,1) (0.5,1) (1,1)};
    \addplot coordinates {(0,1) (0.5,1) (1,1)};
    \addplot coordinates {(0,1) (0.5,1) (1,1)};
    \addplot coordinates {(0,1) (0.5,1) (1,1)};
    \addplot coordinates {(0,1) (0.5,1) (1,1)};
    \addplot coordinates {(0,1) (0.5,1) (1,1)};
    \addplot coordinates {(0,1) (0.5,1) (1,1)};
\end{axis}
\end{tikzpicture}
\end{codeexample}
    %
    此处，我们指定了 $11$ 种颜色并有 $11$ 个绘图。显然，插值颜色的用途
    有限，仅适用于特殊用例。仅当 |colormap| 能适当区分相邻绘图线时，
    才使用此类 |cycle list|！注意，|colors of colormap| 与基于现有
    颜色映射构建颜色映射的相关方式非常相似，如第
    ~\ref{sec:pgfplots:colormaps:based:on:others} 节
    （第~\pageref{sec:pgfplots:colormaps:based:on:others} 页）所述。

    相关选项是 \declareandlabel{indices of colormap}|=|\marg{list}。
    同上，它接受可选的 `|of|' 子句，形式为
    |indices of colormap=|\marg{list}| of |\meta{colormap name}。
    主参数是索引 $0\le N_i<N$ 的列表，其中 $N$ 是颜色映射定义中的
    颜色数量（参见 |index of colormap| 的文档）。超出范围的索引被裁剪到
    最近的索引。例如，|viridis| 有 $\pgfplotscolormapsizeof{viridis}$
    个元素，我们可以写
    %
\begin{codeexample}[]
\begin{tikzpicture}
\begin{axis}[
    colormap name=viridis,colorbar,
    cycle list={
      [indices of colormap={0,4,8,12,17} of viridis]},
    %
    cycle from colormap manual style,
]
    \addplot coordinates {(0,1) (0.5,1) (1,1)};
    \addplot coordinates {(0,1) (0.5,1) (1,1)};
    \addplot coordinates {(0,1) (0.5,1) (1,1)};
    \addplot coordinates {(0,1) (0.5,1) (1,1)};
    \addplot coordinates {(0,1) (0.5,1) (1,1)};
    \addplot coordinates {(0,1) (0.5,1) (1,1)};
    \addplot coordinates {(0,1) (0.5,1) (1,1)};
    \addplot coordinates {(0,1) (0.5,1) (1,1)};
    \addplot coordinates {(0,1) (0.5,1) (1,1)};
    \addplot coordinates {(0,1) (0.5,1) (1,1)};
    \addplot coordinates {(0,1) (0.5,1) (1,1)};
\end{axis}
\end{tikzpicture}
\end{codeexample}
    %
    注意，`| of viridis|' 实际上是多余的，因为在此情况下已选择
    |viridis|。

    有关如何自定义 |of colormap| 的完整语法，与第
    ~\ref{sec:pgfplots:colormaps:based:on:others} 节
    （第~\pageref{sec:pgfplots:colormaps:based:on:others} 页）所述的
    基于其他 |colormap| 定义 |colormap| 的构建块相同——区别在于
    |cycle list={[of colormap]}| 将所选颜色插入 |cycle list| 而非
    |colormap|。

    注意，所有这些特殊列表都是 |\pgfplotscreateplotcyclelist| 的有效参数，
    也可作为 |cycle multi list| 及其变体中的子列表出现。

    由于从 |colormap| 创建 |cycle list| 必然产生没有标记和线型变化的绘图，
    将结果与 |cycle multiindex* list| 组合是有意义的，即连接两个现有列表。
    以下示例将纯颜色列表与标记连接：
    %
\begin{codeexample}[]
\begin{tikzpicture}
\begin{axis}[
    colormap name=hot,colorbar,
    cycle multiindex* list={
        [samples of colormap=11]\nextlist
        mark list\nextlist
    },
    cycle from colormap manual style,
]
    \addplot coordinates {(0,1) (0.5,1) (1,1)};
    \addplot coordinates {(0,1) (0.5,1) (1,1)};
    \addplot coordinates {(0,1) (0.5,1) (1,1)};
    \addplot coordinates {(0,1) (0.5,1) (1,1)};
    \addplot coordinates {(0,1) (0.5,1) (1,1)};
    \addplot coordinates {(0,1) (0.5,1) (1,1)};
    \addplot coordinates {(0,1) (0.5,1) (1,1)};
    \addplot coordinates {(0,1) (0.5,1) (1,1)};
    \addplot coordinates {(0,1) (0.5,1) (1,1)};
    \addplot coordinates {(0,1) (0.5,1) (1,1)};
    \addplot coordinates {(0,1) (0.5,1) (1,1)};
\end{axis}
\end{tikzpicture}
\end{codeexample}
    %
    \noindent 有关 |cycle multiindex* list| 的详情，请参阅下一小节。

    注意，前面的示例都使用以下样式。

    \begin{stylekey}{/pgfplots/cycle from colormap manual style}
        本手册中定义的样式。它的值为
        %
\begin{codeexample}[code only]
\pgfplotsset{
    cycle from colormap manual style/.style={
        x=3cm,y=10pt,ytick=\empty,
        colorbar style={x=,y=,ytick=\empty},
        point meta min=0,point meta max=1,
        stack plots=y,
        y dir=reverse,colorbar style={y dir=reverse},
        every axis plot/.style={line width=2pt},
        legend entries={0,...,20},
        legend pos=outer north east,
    },
}
\end{codeexample}
    \end{stylekey}

}%

    \paragraph{说明：}

    可以在\emph{两个}绘图之间调用 |\pgfplotsset{cycle list=|\marg{a list}|}|
    或 |cycle list name|。此类设置在当前 \TeX{} 组（即花括号）结束前有效。
    每个 |\addplot| 命令使用绘图索引查询 |cycle list|；
    如果期间 |cycle list| 发生变化也不会有问题。
\end{pgfplotskeylist}

\begin{pgfplotskey}{cycle list/.define=\marg{name}\marg{{list}}}
    该命令仅调用
    |\pgfplotscreateplotcyclelist|\allowbreak\marg{name}\allowbreak\marg{list}，
    而不实际将其选为当前列表。

    注意，\PGFPlots{} 使用此命令按以下方式实现其 |cycle list| 键：
    %
\begin{codeexample}[code only]
\pgfplotsset{
    cycle list/.style={%
        cycle list/.define={@internal@}{#1},%
        cycle list name={@internal@}%
    },
}
\end{codeexample}
    %
\end{pgfplotskey}


\subsubsection{组合不同循环列表的构建块}

以下键允许组合不同的 |cycle list| 以构建更复杂的列表。

\begin{pgfplotskey}{%
    cycle multi list=\meta{list 1}\texttt{\textbackslash nextlist}\meta{list 2}%
    \texttt{\textbackslash nextlist}$\ldots$
}
    这是使用多个 |cycle list| 来确定绘图样式的两种方式之一
    （另一种见 |cycle multiindex list|）。最好用示例说明：
    %
\begin{codeexample}[]
\begin{tikzpicture}
    \begin{axis}[
        cycle multi list={
            red,blue\nextlist
            solid,{dotted,mark options={solid}}\nextlist
            mark=*,mark=x,mark=o
        },
        samples=3,
        legend entries={0,...,20},
        legend pos=outer north east,
    ]
        \addplot {x};
        \addplot {x-1};
        \addplot {x-2};
        \addplot {x-3};
        \addplot {x-4};
        \addplot {x-5};
        \addplot {x-6};
        \addplot {x-7};
        \addplot {x-8};
        \addplot {x-9};
        \addplot {x-10};
        \addplot {x-11};
    \end{axis}
\end{tikzpicture}
\end{codeexample}
    %
    \noindent 所提供的 |cycle multi list| 由三个列表组成。单个绘图的样式
    由三个列表中每个列表的元素构成：第一个绘图样式为
    |red,solid,mark=*|，第二个为 |red,solid,mark=x|，第三个为
    |red,solid,mark=o|。第四个绘图重启第三个列表，并使用列表 $2$ 的
    下一个：它具有 |red,dotted,mark options={solid},mark=*|，依此类推。

    最后一个列表始终为新绘图前进。倒数第二个列表（本例中为第二个列表）
    在最后一个列表重置后前进。换言之：|cycle multi list| 允许以字典序
    组合不同的 |cycle list|。\footnote{对于偏好公式的人：索引为
    $0 \le i < N$ 的绘图将使用循环列表偏移 $i_0,i_1,\dotsc,i_k$，
    $0 \le i_m < N_m$，其中 $k$ 是提供给
    \texttt{cycle multi list} 的参数数量，$N_m$ 是第 $m$ 个循环列表
    中的元素数。偏移 $i_m$ 在循环
    {\ttfamily \{ int tmp=i; for( int m=k-1; m>=0; m=m-1 ) \{ i\_m =
    tmp\%N\_m; tmp = tmp/N\_m; \}\}} 中计算。}

    |cycle multi list| 的参数是一系列参数，如为 |cycle list| 提供的那样，
    以 \declareandlabel{\nextlist} 分隔。除提供新循环列表外，
    \meta{list $i$} 元素还可以表示 |cycle list name| 值
    （包括特殊的 |auto| 循环列表，即最近分配的 |cycle list| 或
    |cycle list name|）。最后的 |\nextlist| 是可选的。

    上面示例中的列表也可以写为
    %
\begin{codeexample}[code only]
\begin{axis}[
    cycle multi list={
        red\\blue\\\nextlist
        solid\\dotted,mark options={solid}\\\nextlist
        mark=*\\mark=x\\mark=o\\
    },
]
\end{codeexample}
    %
    \noindent （注意终止符 |\\| 命令！）。

\begin{codeexample}[]
\begin{tikzpicture}
    \begin{axis}[
        title={Cycle color between successive plots, then marks},
        cycle multi list={
            mark list\nextlist
            blue,red
        },
        samples=3,
        legend entries={0,...,20},
        legend pos=outer north east,
    ]
        \addplot {x};
        \addplot {x-1};
        \addplot {x-2};
        \addplot {x-3};
        \addplot {x-4};
        \addplot {x-5};
        \addplot {x-6};
        \addplot {x-7};
        \addplot {x-8};
        \addplot {x-9};
        \addplot {x-10};
        \addplot {x-11};
    \end{axis}
\end{tikzpicture}
\end{codeexample}


    \paragraph{使用子列表}

    \meta{list $i$} 条目也可以使用语法
    |[|\meta{number}| of]|\meta{cycle list name} 仅包含已知循环列表名的
    前 $n$ 个元素。例如 |[2 of]mark list| 将使用 |mark list| 的前 $2$ 个
    元素：
    %
\begin{codeexample}[]
\begin{tikzpicture}
    \begin{axis}[
        title={Cycle 2 marks between successive plots, then colors},
        cycle multi list={
            color list\nextlist
            [2 of]mark list
        },
        samples=3,
        legend entries={0,...,20},
        legend pos=outer north east,
    ]
        \addplot {x};
        \addplot {x-1};
        \addplot {x-2};
        \addplot {x-3};
        \addplot {x-4};
        \addplot {x-5};
        \addplot {x-6};
        \addplot {x-7};
        \addplot {x-8};
        \addplot {x-9};
        \addplot {x-10};
        \addplot {x-11};
    \end{axis}
\end{tikzpicture}
\end{codeexample}
    %
\end{pgfplotskey}

\begin{pgfplotskey}{%
    cycle multiindex list=\meta{list 1}\texttt{\textbackslash nextlist}\meta{list 2}%
    \texttt{\textbackslash nextlist}$\ldots$
}
    这是使用多个 |cycle list| 来确定绘图样式的两种方式之一
    （另一种见 |cycle multi list|）。两种选项的区别在于列表索引如何
    映射到子列表。我们从示例开始：
    %
\begin{codeexample}[]
\begin{tikzpicture}
    \begin{axis}[
        cycle multiindex list={
            red,blue,teal\nextlist
            solid\\dotted,mark options={solid}\\only marks\\\nextlist
            mark=*,mark=x,mark=oplus\nextlist
        },
        samples=3,
        legend entries={0,...,20},
        legend pos=outer north east,
    ]
        \addplot {x};
        \addplot {x-1};
        \addplot {x-2};
        \addplot {x-3};
        \addplot {x-4};
        \addplot {x-5};
        \addplot {x-6};
        \addplot {x-7};
        \addplot {x-8};
        \addplot {x-9};
        \addplot {x-10};
        \addplot {x-11};
    \end{axis}
\end{tikzpicture}
\end{codeexample}
    %
    \noindent 所提供的 |cycle multiindex list| 由三个列表组成。单个绘图的
    样式由三个列表中每个列表的元素构成：第一个绘图样式为
    |red,solid,mark=*|，第二个为
    |blue,dotted,mark options={solid},mark=x|，第三个为
    |teal,only marks,oplus|。第四个绘图重启所有列表，并使用与第一个
    相同的样式，即 |red,solid,mark=*|。

    注意，第二个列表使用列表分隔符 `|\\|'，它需要为 |cycle list|
    定义的最终终止符。

    因此，该样式对每个列表使用\emph{相同}索引（``多索引''）。
    因此，它的不同选项比 |cycle multi list| 少得多（后者产生所有可能的
    变化），但其组合方法解决不同的用例。

    |cycle multiindex list| 的参数格式与
    与 |cycle multi list| 相同，包括特殊的 |[2 of]mark list| 语法
    以及按名称提供其他周期列表：
    %
\begin{codeexample}[]
\begin{tikzpicture}
    \begin{axis}[
        title={Same index of each input list},
        cycle multiindex list={
            [3 of]mark list\nextlist
            blue,red\nextlist
            linestyles\nextlist
        },
        samples=3,
        legend entries={0,...,20},
        legend pos=outer north east,
    ]
        \addplot {x};
        \addplot {x-1};
        \addplot {x-2};
        \addplot {x-3};
        \addplot {x-4};
        \addplot {x-5};
        \addplot {x-6};
        \addplot {x-7};
        \addplot {x-8};
        \addplot {x-9};
        \addplot {x-10};
        \addplot {x-11};
    \end{axis}
\end{tikzpicture}
\end{codeexample}

    注意，|cycle multiindex list| 接受不同大小的列表。|cycle multiindex list|
    的大小是\emph{最大}输入列表的大小，所有较小的输入列表用空选项列表填充。
    这就是为什么前一个示例每第三个绘图使用颜色 |black|：第二个列表中
    没有颜色，省略颜色导致 |black|。一旦最大子列表的最后一项被使用，
    列表就重启。
\end{pgfplotskey}

\begin{pgfplotskey}{%
    cycle multiindex* list=\meta{list 1}\texttt{\textbackslash nextlist}\meta{list 2}%
    \texttt{\textbackslash nextlist}$\ldots$
}
    |cycle multiindex list| 的一个变体，行为相同——除了不同大小的子列表。

    如上所述，非星号版本 |cycle multiindex list| 用\emph{空}选项列表填充
    缺失条目，直到所有列表元素大小相同。

    星号键 |cycle multiindex* list| 在子列表到达末尾时独立重启：
    %
\begin{codeexample}[]
\begin{tikzpicture}
    \begin{axis}[
        title={Same index of each input list, lists restarted},
        cycle multiindex* list={
            [3 of]mark list\nextlist
            blue,red\nextlist
            linestyles\nextlist
        },
        samples=3,
        legend entries={0,...,20},
        legend pos=outer north east,
    ]
        \addplot {x};
        \addplot {x-1};
        \addplot {x-2};
        \addplot {x-3};
        \addplot {x-4};
        \addplot {x-5};
        \addplot {x-6};
        \addplot {x-7};
        \addplot {x-8};
        \addplot {x-9};
        \addplot {x-10};
        \addplot {x-11};
    \end{axis}
\end{tikzpicture}
\end{codeexample}
    %
    这与非星号变体 |cycle multiindex list| 文档中记录的示例完全相同。
    但是，第二个子列表元素较少——非星号变体产生 |black|，而星号变体
    在两个现有颜色用完后重启第二个子列表。

    该样式允许以复杂方式连接列表：
    %
\begin{codeexample}[]
\begin{tikzpicture}
    \begin{axis}[
        title={Same index of each input list},
        cycle multiindex* list={
            red,blue,teal,brown\nextlist
            mark=*,mark=square*,mark=triangle*\nextlist
            dashed,every mark/.append style={
                solid,draw=.!50!black,fill=.,}\\
            \nextlist
        },
        samples=3,
        legend entries={0,...,20},
        legend pos=outer north east,
    ]
        \addplot {x};
        \addplot {x-1};
        \addplot {x-2};
        \addplot {x-3};
        \addplot {x-4};
        \addplot {x-5};
        \addplot {x-6};
        \addplot {x-7};
        \addplot {x-8};
        \addplot {x-9};
        \addplot {x-10};
        \addplot {x-11};
    \end{axis}
\end{tikzpicture}
\end{codeexample}
    %
    我们看到四种不同颜色按预期周期性出现。我们还看到三种不同标记
    有各自的周期（如预期每四个绘图重启）。但第三个子列表只包含一个
    元素，我们可以通过其分隔符 `|\\|' 看到，它只在列表末尾出现一次！
    因此，该列表为\emph{每个}绘图重启，使每个绘图都获得其参数。
\end{pgfplotskey}


\subsection{坐标轴背景}

\begin{pgfplotskey}{axis background (initially empty)}
    这是用于配置坐标轴本身外观的样式。可以使用
    |axis background/.style=|\marg{options} 方法定义和/或修改它。
    将使用 \meta{options} 生成背景路径，其中可包含填充颜色或底纹。
    %
\pgfplotsexpensiveexample
\begin{codeexample}[]
\begin{tikzpicture}
\begin{axis}[
    axis background/.style={fill=blue!10},
]
    \addplot3 [
        surf,
        y domain=0:1,
    ]
        {sin(deg(x)) * y*(1-y)};
\end{axis}
\end{tikzpicture}
\end{codeexample}

    请注意，在默认配置中图例填充为白色。
    %
\begin{codeexample}[]
\begin{tikzpicture}
\begin{semilogyaxis}[
    axis background/.style={
        shade,top color=gray,bottom color=white,
    },
    legend style={fill=white},
]
    \addplot {exp(-x)};
    \addplot {exp(-4*x)};

    \legend{$e^{-x}$,$e^{-4x}$}
\end{semilogyaxis}
\end{tikzpicture}
\end{codeexample}
    %
    有关 |fill| 和 |shade| 的详情，请参阅 \Tikz{} 手册
    \cite{tikz}。
\end{pgfplotskey}


\section{提供颜色数据——元数据}
\label{sec:pgfplots:point:meta}

\PGFPlots{} 提供根据特殊坐标（``元数据''）修改绘图的功能。例如，
散点图可能根据此特殊数据改变标记颜色、大小或外观。曲面图和网格图是
另一个示例：此处，曲面补丁（或网格部分）的颜色取决于``元数据''。
节点的``元数据''不是位置（位置由 $(x,y)$ 或 $(x,y,z)$ 给出）。
它是关于该节点的``某些信息''；可以说：它是节点的特定属性。
这通常称为``元数据''。


\subsection{元数据概述}

在 \PGFPlots{} 中，每个节点都有坐标\emph{和}元数据。因此，二维绘图有
三个值：$x$、$y$ 和 |point meta| 数据。三维绘图每个坐标有四个值：
$x$、$y$、$z$ 和 |point meta|。

在许多情况下，|point meta| 被解释为颜色数据。更准确地说：它被解释为
\emph{标量}颜色数据，映射到 |colormap|（下一段将详细介绍此方法）。
但是，|point meta| 可以是任何内容。通常，它是一个数字，如颜色数据的情况。
该数字线性映射到区间 $[0,1000]$，使 $0$ 对应遇到的最小值，$1000$
对应遇到的最大标量值。映射后的值可作为
|\pgfplotspointmetatransformed| 使用。此特殊值允许定义绘图的某些属性：
它可以是颜色（配合 |colormap|）。它也可以是 |mesh| 图中的
|line width|（更准确地说：|line width| 可定义为取决于变换后的元数据）。
此方法的示例见第~\pageref{pgfplots:example:pointmeta:quiver} 页，
其中 |quiver| 图的箭头线取决于 |\pgfplotspointmetatransformed|。
该值也可用于定义 |mark size|，如
|\scatter/@pre marker code| 在第~\pageref{pgfplots:example:pointmeta:scatter}
页的示例。但是，|point meta| 数据不一定是数字。它可以是文本标签
（任何你喜欢的文本）。例如，|nodes near coords| 使用它。它也可以包含
元组，如 RGB 颜色信息（参见 |mesh/color input=explicit|）。
因此，|point meta| 确实是关于单个坐标的某些抽象信息。

注意，每个点只有\emph{一个} |point meta|。如果你需要每个坐标多个元数据
值，请参见键 |visualization depends on|。

共同思路是告诉 \PGFPlots{} 如何获取元数据。无需显式提供数据——在许多
情况下，用于着色 |surf|ace 补丁或 |mark|er 颜色的数据是绘图的 $y$ 或
$z$ 坐标。告诉 \PGFPlots{} 在何处找到``元数据''的方法是
|point meta| 键。

|point meta| 最常见的用例是颜色信息：如果元数据在区间
$[m_{\min},m_{\max}]$ 内，遇到的最小值 $m=m_{\min}$ 的元数据坐标将获得
|colormap| 提供的最低颜色，而遇到的最大值 $m=m_{\max}$ 将获得
|colormap| 提供的最高颜色。如前所述，这通过
|\pgfplotspointmetatransformed| $\in [0,1000]$（按惯例）完成。
最小和最大值之间的任何坐标都将线性映射：例如，均值
$m=1/2(m_{\max}+m_{\min})$ 将获得颜色映射的中间颜色
（其 |\pgfplotspointmetatransformed| $=500$）。这就是为什么本手册中
``元数据''有时称为``颜色数据''。

\begin{codeexample}[]
\begin{tikzpicture}
    \begin{axis}[colorbar]
        \addplot [mesh,point meta=y,thick] {x^2};
    \end{axis}
\end{tikzpicture}
\end{codeexample}


\subsection{元数据的用户输入格式}

\begin{pgfplotskey}{point meta=\mchoice{%
        none,
        \meta{expression},
        x,y,z,f(x),
        explicit,explicit symbolic%
    } (initially none)%
}
\label{pgfplots:pointmeta}
    |point meta| 键告诉 \PGFPlots{} 在何处获取特殊元数据。请注意，
    |point meta| 和 |scatter src| 实际上是相同的——|scatter src| 是
    |point meta| 的别名。因此，第~\pageref{pgfplots:scatter:src} 页
    为 |scatter src| 提供的摘要涵盖相同主题。不过，|point meta| 的主要
    参考在这里。
    %
    \begin{description}
        \item[\declaretext{none}] 初始选项 |none| 禁用元数据，
            不产生计算开销。任何其他选项将激活元数据上下限范围的计算，
            即 $[m_{\min},m_{\max}]$ 的计算。
        \item[\declaretext{x}] 选项 |x| 使用已有的 $x$ 坐标作为元数据。
            这始终指应用任何用户变换、对数、堆叠绘图计算等之后的
            \emph{最终} $x$ 坐标。如果你需要此处的未处理坐标值，
            请考虑使用 |rawx|。
        \item[\declaretext{y}]
        \item[\declaretext{z}] 选项 |y| 和 |z| 类似：它们分别使用
            $y$ 或 $z$ 坐标作为元数据。因此，这三个选项\emph{不需要}
            任何额外数据。与 |x| 一样，有数学常量 |rawy| 和 |rawz|，
            分别产生未处理的 $y$ 和 $z$ 值。
        \item[\declaretext{f(x)}] 这将使用最后一个可用坐标，
            换言之：二维绘图中与 |y| 相同，三维绘图中与 |z| 相同。
        \item[\declaretext{explicit}] 该选项告诉 \PGFPlots{} 期望在
            坐标输入流中显式提供\emph{数值}元数据。此数据将线性变换到
            当前颜色映射，如前所述。

            如何为 |\addplot coordinates|、|\addplot table| 和其他输入方法
            提供元数据，在第~\ref{sec:pgfplots:providing:input} 节中详细
            描述——但我们在此提供小示例来总结各种可能性：
            %
\begin{codeexample}[code only]
% for 'coordinates':
% provide color data explicitly using [<data>]
% behind coordinates:
\addplot+ [point meta=explicit]
    coordinates {
        (0,0) [1.0e10]
        (1,2) [1.1e10]
        (2,3) [1.2e10]
        (3,4) [1.3e10]
        % ...
    };
\end{codeexample}

\begin{codeexample}[code only]
% for 'table':
% Assumes a datafile.dat like
% xcolname  ycolname    colordata
% 0         0           0.001
% 1         2           0.3
% 2         2.1         0.4
% 3         3           0.5
% ...
% the file may have more columns.
\addplot+ [point meta=explicit]
    table [x=xcolname,y=ycolname,meta=colordata]
        {datafile.dat};
% or, equivalently (perhaps a little bit slower):
\addplot+ [point meta=\thisrow{colordata}]
    table [x=xcolname,y=ycolname]
        {datafile.dat};
\end{codeexample}

\begin{codeexample}[code only]
% for 'file':
% Assumes a datafile.dat like
% 0         0           0.001
% 1         2           0.3
% 2         2.1         0.4
% 3         3           0.5
% ...
% the first three columns will be used here as x,y and meta,
% resp.
\addplot+ [point meta=explicit]
    file {datafile.dat};
\end{codeexample}

\begin{codeexample}[code only]
% 'table' using expressions which may depend on all
% columns:
% Assumes a datafile.dat like
% xcolname  ycolname    anything    othercol
% 0         0           4           15
% 1         2           5           20
% 2         2.1         8           30
% 3         3           42          40
% ...
% the file may have more columns.
\addplot+ [point meta={0.5*(\thisrow{anything} + sqrt(\thisrow{othercol}))}]
    table [x=xcolname,y=ycolname]
        {datafile.dat};
\end{codeexample}
            %
            因此，有多种方法提供元数据（颜色数据）。选项 |explicit| 的
            关键在于显式提供某些数据——尽管 |point meta| 不知道如何提供。
            数据应为数值类型，并线性映射到范围 $[0,1000]$
            （可能用于当前颜色映射）。
        \item[\declaretext{explicit symbolic}] 选项 |explicit symbolic|
            与 |explicit| 非常相似，它期望坐标输入例程提供额外数据。
            但是，|explicit symbolic| 不一定期望数值数据：你可以提供
            任何类型的符号。可以提供一组样式，散点图中每个类一个。
            这使用 |scatter/classes| 实现，见第
            ~\pageref{pgfplots:scatterclasses} 页。输入数据的提供方式与
            上述选项 |explicit| 相同。

            该选项适用于带文本标签的 |nodes near coords|（见第
            ~\pageref{pgfplots:example:pointmeta:nodesnearcoords} 页）
            和带显式颜色的 |surf|ace 图（详情见
            |mesh/color input=explicit|）。
        \item[\normalfont\declare{\meta{expression}}] 该选项允许使用数学
            表达式计算元数据。\meta{expression} 可取决于 |x|、|y|、|z|，
            分别产生当前 $x$、$y$ 或 $z$ 坐标。坐标已完全处理
            （变换、对数），如上述选项 |x| 所述。此外，\meta{expression}
            可取决于 |\addplot| 期间有效的命令，如 |\plotnum| 或
            |\coordindex|（详见第~\ref{sec:pgfplots:misc} 节）。如果使用
            |\addplot table| 提供坐标，宏 |\thisrow|\marg{colname}
            特别有用，因为它允许访问``当前''行中提供的
            \meta{colname} 的值。计算使用 \PGF{} 的浮点单元执行，
            所有支持的算术运算均可使用。

            本质上，\meta{expression} 可取决于所有 |\addplot| 命令已知的
            一切：$x$、$y$ 和（如果有）$z$ 坐标。此外，它可取决于
            |rawx|、|rawy| 或 |rawz|。这三个表达式产生输入流中找到的
            未处理 $x$、$y$ 或 $z$ 值（无对数、无用户变换）。%
            \footnote{%
                在极少数情况下，将数学表达式应用于元数据的另一个源
                （其他选项之一）可能很有趣。为此，在其他可能选项已求值后，
                才检查 \meta{expression}。换言之，语句
                \texttt{point meta=explicit,
                point meta=meta*meta+3} 将使用设为显式提供数据的
                |meta| 来求值表达式。%
            }
            如果与 |\addplot table| 一起使用，你还可以访问其他表格列
            （例如使用 |\thisrow|\marg{colname}）。
        \item[\normalfont\declaretext{TeX code}\texttt{=}\meta{code}] 一个
            相当底层的选项，允许提供 \TeX{} \meta{code} 来计算数值。
            \meta{code} 应定义宏 |\pgfplotspointmeta|。它在局部作用域
            环境中求值（其局部变量之后被释放）。它可取决于与上述
            \meta{expression} 相同的值，特别是 |\thisrow|\marg{colname}
            for table input.

            注意，此时数学解析器将配置为使用 |fpu|，因此 |\pgfmathparse|
            产生浮点数。

            注意，需要额外一对花括号来提供此键，即

            |point meta={TeX code=|\marg{code}|}|.

        \item[\normalfont\declaretext{TeX code
            symbolic}\texttt{=}\meta{code}] 与 |TeX code| 一样，可提供
            定义宏 |\pgfplotspointmeta| 的 \meta{code}，但结果不解释为数字。
            它类似 |explicit symbolic| 选项。

            注意，需要额外一对花括号来提供此键，即

            |point meta={TeX code symbolic=|\marg{code}|}|.
        \item[\normalfont\declaretext{symbolic}\texttt{=}\meta{symbol}] 一个
            选项，接受任意 \meta{symbol}，用于每个坐标。与
            |explicit symbolic| 和 |TeX code symbolic| 一样，此选项产生
            符号表示，即它保持原样，不映射结果。

            与 |explicit symbolic| 的区别是，\meta{symbol} 是所有点的公共
            符号，而 |explicit symbolic| 期望输入数据流（如表格）提供
            单独的符号。

            与 |TeX code symbolic| 的区别很小：|symbolic| 实际上等同于

            |point meta/TeX code symbolic={\def\pgfplotspointmeta|\marg{symbol}|}|.

            |symbolic| 的一个用例是 |mesh/color input=explicit mathparse|，
            见其中文档。

            注意，需要额外一对花括号来提供此键，即

            |point meta={symbolic=|\marg{symbol}|}|.
    \end{description}

    如前所述，元数据的主要应用是使用到范围 $[0,1000]$ 的线性映射
    （可能用于当前颜色映射）来确定（标记/面/边）颜色。此映射工作如下：
    它是函数
        \[ \phi\colon [m_{\min},m_{\max}] \to [0,1000] \]
    其中
        \[ \phi(m) = \frac{m - m_{\min}}{m_{\max} - m_{\min}} \cdot 1000 \]
    使 $\phi(m_{\min})=0$ 且 $\phi(m_{\max})=1000$。按惯例，
    值 $1000$ 是所有颜色映射的上限。现在，如果一个坐标（或边/面）
    具有元数据 $m$，其颜色将使用 $\phi(m)$ 确定：它是当前颜色映射中
    $\phi(m)$\textperthousand{} 处的颜色。

    此变换取决于区间 $[m_{\min},m_{\max}]$，而该区间又可使用下文描述的
    键 |point meta rel|、|point meta min| 和 |point meta max| 修改。

    未变换的元数据可在宏
    \declareandlabel{\pgfplotspointmeta} 中获取（仅在正确上下文中，
    例如散点图样式或 |scatter/@pre marker code| 接口）。此宏包含
    底层浮点数（除非它是非解析字符串数据）。变换后的数据可在宏
    \declareandlabel{\pgfplotspointmetatransformed} 中获取，为定点表示。
    预期其在范围 $[0,1000]$ 内。

\end{pgfplotskey}

\begin{pgfplotskey}{set point meta if empty=\marg{point meta source}}
    设置 |point meta=|\meta{point meta source}，但仅在当前
    |point meta=none| 时。这用于 |scatter|、|mesh| 和 |surf|，
    配合 |set point meta if empty=f(x)|。
\end{pgfplotskey}


\subsection{映射元数据与颜色映射}

如 |point meta| 文档中已解释的，元数据的一个应用是使用当前颜色映射
和从元数据到当前颜色映射的线性映射来确定颜色。本小节显示控制此映射的键。

绘图处理程序 |surf|、|mesh|、|patch|、|scatter|、|contour prepared|
及一些相关处理程序都使用映射颜色。为此，|point meta| 提供的值按以下
方式处理：
%
\begin{itemize}
    \item 首先，收集所有遇到的值（在 |\addplot| 期间），并计算最大值和
        最小值 $m_{\max}$ 和 $m_{\min}$。值存储在 |\pgfplotspointmeta| 中。
    \item 其次，使用线性映射将每个遇到的值 $m$ 映射到 $[0,1000]$：
        $m_{\min}\to 0$；$m_{\max}\to 1000$，中间所有值线性映射。
        此值存储（并可获取）在 \TeX{} 宏 |\pgfplotspointmetatransformed|
        中（但在 |\end{axis}| 之前不可用）。
    \item 最后，映射值 |\pgfplotspointmetatransformed| 被映射到当前
        |colormap|。生成的颜色称为
        \declareandlabel{mapped color}。
\end{itemize}
%
绘图处理程序进而使用 |mapped color| 来渲染颜色
（通常与颜色混合组合，如 |mapped color!80!black|）。

\begin{pgfplotskey}{point meta rel=\mchoice{axis wide,per plot} (initially axis wide)}
    键 |point meta rel| 配置所有元数据坐标的区间 $[m_{\min},m_{\max}]$
    是作为整个坐标轴中所有绘图的最大值计算（选项
    \declaretext{axis wide}），还是仅针对单个特定绘图
    （选项 \declaretext{per plot}）。

\message{Overfull hbox is ok.}%
\begin{codeexample}[]
\begin{tikzpicture}
    \begin{axis}[
        title=Axis wide color mapping,
        colorbar,
        samples=50,point meta rel=axis wide,
        point meta=y,
    ]
        \addplot [mesh,thick] {sin(deg(x))};
        \addplot [mesh,thick] {3*tanh(x)};
    \end{axis}
\end{tikzpicture}
~
\begin{tikzpicture}
    \begin{axis}[
        title=Per Plot color mapping,
        colorbar,
        samples=50,
        point meta rel=per plot,
        point meta=y,
    ]
        \addplot [mesh,thick] {sin(deg(x))};
        \addplot [mesh,thick] {3*tanh(x)};
   \end{axis}
\end{tikzpicture}
\end{codeexample}

    注意，|colorbar| 仍将使用 |axis wide| 元数据限制。如果你希望颜色条使用
    某个\emph{特定}绘图的颜色数据限制，请考虑 |colorbar source| 键。
    |point meta rel| 键配置元数据如何映射到 |colormap| 中的颜色。
\end{pgfplotskey}

\begin{pgfplotskeylist}{point meta min=\marg{number},point meta max=\marg{number}}
    这些键允许显式定义元数据到范围 $[0,1000]$ 的线性映射所需的范围
    （例如，用于当前映射）。如果同一映射要用于多个坐标轴，则这是必要的。


    \paragraph{特殊情况说明：}

    \begin{itemize}
        \item 可以部分提供限制；此时仅计算缺失的限制。
        \item 如果元数据落在这些限制之外，线性变换仍有良好定义，
            这是可接受的（除非区间长度为零）。但是，颜色数据不能
            落在这些限制之外，因此颜色条会执行截断。
        \item 该键既可用于单个绘图，也可用于整个坐标轴（或两者）。
        \item 如果为单个绘图提供了元数据限制，这些限制也可能贡献到
            坐标轴范围元数据区间。
    \end{itemize}
\end{pgfplotskeylist}

\begin{pgfplotskey}{colormap access=\mchoice{%
        map,
        direct,
        const,\\
        piecewise constant,
        piecewise const,
        piecewise linear%
    } (initially map)%
}
    该键配置元数据如何用于从颜色映射确定颜色。


    \subsubsection{使用分段常数插值映射元数据}

    初始配置 \declaretext{map} 执行上面详细解释的线性映射操作。


    \subsubsection{将颜色映射用作索引颜色空间（调色板）}

    选项 \declaretext{direct} 不执行任何变换；它将 |point meta| 的值
    解释为当前颜色映射的整数索引。如果数字不是整数，将通过丢弃小数部分
    来截断。因此，它与 |index of colormap| 非常相似。\index{index of colormap}

\def\showcolorandname#1{%
    \showcolor{#1}~\pgfmanualpdfref{#1}{\texttt{#1}}%
}%
\def\showcolor#1{%
    \tikz \draw[black,fill={#1}] (0,0) rectangle (1em,0.6em);%
}%
%
    假设我们有颜色定义
    %
    \showcolorandname{RdPu-K},
    \showcolorandname{RdPu-I},
    \showcolorandname{RdPu-G},
    \showcolorandname{RdPu-E}, 和
    \showcolorandname{RdPu-B}
    %
    （实际上取自 |\usetikzlibrary{colorbrewer}|）。

    我们将使用这些颜色定义一个 |colormap|：
    %
\begin{codeexample}[]
% requires \usetikzlibrary{colorbrewer}
\begin{tikzpicture}
\begin{axis}[
    colormap={brewer}{
        color=(RdPu-B) color=(RdPu-E)
        color=(RdPu-G) color=(RdPu-I)
        color=(RdPu-K)
    },
    enlargelimits=false,
    colorbar as palette,
]
    \addplot [
        matrix plot,
        nodes near coords,mark=*,
        point meta=explicit,
        colormap access=direct,
    ] coordinates {
        (0,0) [0] (1,0) [1] (2,0) [2] (3,0) [NaN]

        (0,1) [3] (1,1) [4] (2,1) [5] (3,1) [99]

        (0,2) [3.5] (1,2) [-1] (2,2) [2.99] (3,2) [3]
    };
\end{axis}
\end{tikzpicture}
\end{codeexample}
    %
    该示例有一个 $3\times 4$ 矩阵，每个都有显式提供的 |point meta| 值。
    第一个点的 |point meta| 值为~$0$，使用 |colormap| 的第一个颜色
    （|RdPu-B|，\showcolor{RdPu-B}）。第二个点值为~$1$，使用第二个颜色
    |RdPu-E| \showcolor{RdPu-E}，依此类推。注意，|NaN| 被解释为矩阵中的
    ``空洞''，留空（此处你看到背景色）。此外，注意坐标
    |(1,2) [-1]| 被映射到第一个颜色，因为超出允许范围的值会自动裁剪到
    下一个正确颜色。点 |[99]| 同理：它被映射到最后一个可用颜色
    ~\showcolor{RdPu-K}。最后，点 |[2.99]| 被映射到索引 $2$
    ~(\showcolor{RdPu-G})，因为 |colormap access=direct| 截断到下一个
    整数索引。

    \begin{pgfplotskey}{colorbar as palette}
        上面的示例使用了 |colorbar as palette|，这是一种专用样式，
        适用于此类调色板型应用。

        它定义为
        %
\begin{codeexample}[code only]
\pgfplotsset{
    colorbar as palette/.style={
        colorbar sampled={
            surf,
            colormap access=direct,
            shader=flat corner,
            samples=\pgfplotscolormapsizeof
            {\pgfkeysvalueof{/pgfplots/colormap name}}+1,
        },
        colorbar style={
            point meta min=0,
            point meta max=\pgfplotscolormapsizeof
            {\pgfkeysvalueof{/pgfplots/colormap name}},
        },
    },
}
\end{codeexample}
        %
        其思路是仅显示``调色板''的颜色，即不插值 |colormap| 的值。
        此外，它应该显示所有可用颜色。

        另见密切相关（但更通用）的 |colorbar as legend|。
\end{pgfplotskey}

    注意，|colormap access=direct| 改变 |point meta| 转换为颜色的方式。
    但是，它\emph{不影响} |shader|。换言之，使用 |shader=interp| 或
    |shader=flat mean| 完全合理。但是，这将插值由
    |colormap access=direct| 产生的颜色！我们上面的示例使用
    |shader=flat corner|，完全不进行颜色插值。

    选项 |colormap access=direct| 实际上就像带索引颜色空间的图像的颜色
    调色板：假设我们使用随机选择的图像，将其降采样使其对 \LaTeX{}
    不致过大，并导出为以下形式的坐标：
    %
\begin{codeexample}[code only]
1 1 194
1 2 194
1 3 170
1 4 193
1 5 188
1 6 188
...
\end{codeexample}
    %
    带索引颜色空间。此外，我们将其颜色调色板导入为 \PGFPlots{}
    |colormap|。为此，一个小的 Matlab 脚本用作将 |tiff| 图像转换为
    坐标的工具：\footnote{另见第~\ref{sec:pgfplots:import:matlab} 节，
    ``从 Matlab 导入网格数据到 PGFPlots''，第
    ~\pageref{sec:pgfplots:import:matlab} 页}
    %
\begin{codeexample}[code only]
[Z,cm] = imread('legoplane.tiff');
size(Z)
[X,Y]=meshgrid(1:size(Z,2), 1:size(Z,1));
data = [ X(:) Y(:) Z(:) ];
fid = fopen('legoplane.dat', 'w');
fprintf(fid, "%d %d %d\n", data');
fclose(fid);
fid = fopen('legoplane_cm.tex', 'w');
fprintf(fid, "\\pgfplotsset{colormap={legoplane palette}{\n");
fprintf(fid, "rgb=(%f,%f,%f)\n", cm');
fprintf(fid, "}}\n");
fclose(fid);
\end{codeexample}
    %
    这允许我们将图形导入 \PGFPlots{} 作为 |matrix plot|：
    %
\tikzset{external/figure name/.add={}{legoplane_}}%
\pgfplotsexpensiveexample
\begin{codeexample}[]
\begin{tikzpicture}
    % load and activate the CM:
    \pgfplotsset{colormap={legoplane palette}{
rgb=(0.058824,0.043137,0.082353)
rgb=(0.078431,0.047059,0.070588)
rgb=(0.054902,0.050980,0.105882)
rgb=(0.098039,0.054902,0.101961)
rgb=(0.109804,0.054902,0.082353)
rgb=(0.172549,0.035294,0.098039)
rgb=(0.129412,0.062745,0.074510)
rgb=(0.070588,0.074510,0.137255)
rgb=(0.054902,0.078431,0.164706)
rgb=(0.027451,0.086275,0.243137)
rgb=(0.113725,0.078431,0.074510)
rgb=(0.027451,0.105882,0.109804)
rgb=(0.141176,0.074510,0.105882)
rgb=(0.125490,0.082353,0.105882)
rgb=(0.105882,0.086275,0.129412)
rgb=(0.250980,0.050980,0.105882)
rgb=(0.192157,0.070588,0.090196)
rgb=(0.211765,0.078431,0.082353)
rgb=(0.223529,0.078431,0.066667)
rgb=(0.188235,0.086275,0.101961)
rgb=(0.164706,0.094118,0.101961)
rgb=(0.207843,0.098039,0.070588)
rgb=(0.192157,0.101961,0.086275)
rgb=(0.403922,0.050980,0.090196)
rgb=(0.462745,0.035294,0.086275)
rgb=(0.254902,0.101961,0.082353)
rgb=(0.290196,0.094118,0.090196)
rgb=(0.380392,0.066667,0.129412)
rgb=(0.180392,0.125490,0.152941)
rgb=(0.211765,0.125490,0.105882)
rgb=(0.160784,0.133333,0.180392)
rgb=(0.239216,0.117647,0.113725)
rgb=(0.050980,0.149020,0.368627)
rgb=(0.235294,0.121569,0.133333)
rgb=(0.207843,0.133333,0.141176)
rgb=(0.239216,0.129412,0.094118)
rgb=(0.184314,0.149020,0.082353)
rgb=(0.094118,0.152941,0.317647)
rgb=(0.313725,0.117647,0.078431)
rgb=(0.015686,0.196078,0.180392)
rgb=(0.050980,0.192157,0.156863)
rgb=(0.529412,0.058824,0.082353)
rgb=(0.141176,0.156863,0.258824)
rgb=(0.247059,0.133333,0.184314)
rgb=(0.047059,0.172549,0.435294)
rgb=(0.160784,0.160784,0.231373)
rgb=(0.596078,0.050980,0.105882)
rgb=(0.301961,0.141176,0.086275)
rgb=(0.200000,0.176471,0.074510)
rgb=(0.227451,0.156863,0.200000)
rgb=(0.349020,0.129412,0.180392)
rgb=(0.215686,0.168627,0.192157)
rgb=(0.337255,0.141176,0.125490)
rgb=(0.549020,0.078431,0.129412)
rgb=(0.274510,0.156863,0.172549)
rgb=(0.196078,0.176471,0.223529)
rgb=(0.309804,0.152941,0.129412)
rgb=(0.282353,0.160784,0.149020)
rgb=(0.282353,0.164706,0.133333)
rgb=(0.505882,0.098039,0.160784)
rgb=(0.356863,0.149020,0.101961)
rgb=(0.258824,0.172549,0.168627)
rgb=(0.376471,0.152941,0.082353)
rgb=(0.392157,0.156863,0.070588)
rgb=(0.321569,0.168627,0.243137)
rgb=(0.086275,0.223529,0.486275)
rgb=(0.015686,0.243137,0.552941)
rgb=(0.352941,0.188235,0.164706)
rgb=(0.188235,0.223529,0.333333)
rgb=(0.662745,0.105882,0.117647)
rgb=(0.247059,0.215686,0.262745)
rgb=(0.658824,0.105882,0.141176)
rgb=(0.211765,0.223529,0.301961)
rgb=(0.274510,0.211765,0.239216)
rgb=(0.364706,0.196078,0.133333)
rgb=(0.333333,0.200000,0.192157)
rgb=(0.654902,0.109804,0.160784)
rgb=(0.321569,0.203922,0.215686)
rgb=(0.603922,0.125490,0.176471)
rgb=(0.317647,0.219608,0.117647)
rgb=(0.298039,0.215686,0.219608)
rgb=(0.654902,0.113725,0.184314)
rgb=(0.329412,0.211765,0.172549)
rgb=(0.309804,0.215686,0.192157)
rgb=(0.509804,0.160784,0.160784)
rgb=(0.160784,0.239216,0.447059)
rgb=(0.486275,0.164706,0.227451)
rgb=(0.411765,0.196078,0.145098)
rgb=(0.431373,0.203922,0.105882)
rgb=(0.431373,0.203922,0.129412)
rgb=(0.215686,0.262745,0.231373)
rgb=(0.415686,0.207843,0.192157)
rgb=(0.239216,0.250980,0.290196)
rgb=(0.403922,0.231373,0.109804)
rgb=(0.639216,0.164706,0.094118)
rgb=(0.439216,0.211765,0.254902)
rgb=(0.419608,0.215686,0.298039)
rgb=(0.419608,0.235294,0.137255)
rgb=(0.290196,0.274510,0.141176)
rgb=(0.403922,0.235294,0.203922)
rgb=(0.419608,0.235294,0.168627)
rgb=(0.258824,0.270588,0.352941)
rgb=(0.321569,0.258824,0.286275)
rgb=(0.376471,0.247059,0.250980)
rgb=(0.403922,0.247059,0.172549)
rgb=(0.349020,0.250980,0.298039)
rgb=(0.294118,0.266667,0.317647)
rgb=(0.619608,0.176471,0.254902)
rgb=(0.376471,0.254902,0.207843)
rgb=(0.349020,0.258824,0.250980)
rgb=(0.235294,0.278431,0.396078)
rgb=(0.454902,0.235294,0.184314)
rgb=(0.035294,0.329412,0.639216)
rgb=(0.678431,0.196078,0.066667)
rgb=(0.537255,0.250980,0.062745)
rgb=(0.466667,0.274510,0.211765)
rgb=(0.447059,0.294118,0.098039)
rgb=(0.466667,0.282353,0.160784)
rgb=(0.466667,0.282353,0.180392)
rgb=(0.262745,0.317647,0.431373)
rgb=(0.513725,0.270588,0.168627)
rgb=(0.576471,0.239216,0.298039)
rgb=(0.450980,0.286275,0.207843)
rgb=(0.466667,0.278431,0.247059)
rgb=(0.450980,0.278431,0.298039)
rgb=(0.415686,0.294118,0.250980)
rgb=(0.376471,0.298039,0.345098)
rgb=(0.411765,0.294118,0.286275)
rgb=(0.349020,0.301961,0.396078)
rgb=(0.352941,0.305882,0.356863)
rgb=(0.215686,0.329412,0.529412)
rgb=(0.498039,0.278431,0.223529)
rgb=(0.450980,0.290196,0.254902)
rgb=(0.321569,0.317647,0.364706)
rgb=(0.301961,0.321569,0.403922)
rgb=(0.301961,0.349020,0.176471)
rgb=(0.447059,0.294118,0.364706)
rgb=(0.396078,0.321569,0.262745)
rgb=(0.505882,0.290196,0.274510)
rgb=(0.388235,0.333333,0.231373)
rgb=(0.407843,0.345098,0.082353)
rgb=(0.509804,0.325490,0.043137)
rgb=(0.533333,0.301961,0.215686)
rgb=(0.309804,0.356863,0.392157)
rgb=(0.400000,0.337255,0.329412)
rgb=(0.533333,0.309804,0.250980)
rgb=(0.505882,0.317647,0.286275)
rgb=(0.517647,0.317647,0.254902)
rgb=(0.505882,0.325490,0.219608)
rgb=(0.698039,0.258824,0.317647)
rgb=(0.501961,0.329412,0.196078)
rgb=(0.419608,0.337255,0.380392)
rgb=(0.498039,0.321569,0.325490)
rgb=(0.490196,0.329412,0.290196)
rgb=(0.498039,0.333333,0.247059)
rgb=(0.376471,0.352941,0.447059)
rgb=(0.466667,0.337255,0.337255)
rgb=(0.364706,0.360784,0.411765)
rgb=(0.313725,0.368627,0.490196)
rgb=(0.400000,0.352941,0.407843)
rgb=(0.337255,0.368627,0.450980)
rgb=(0.447059,0.356863,0.247059)
rgb=(0.403922,0.364706,0.384314)
rgb=(0.698039,0.286275,0.372549)
rgb=(0.192157,0.411765,0.674510)
rgb=(0.592157,0.325490,0.376471)
rgb=(0.513725,0.364706,0.278431)
rgb=(0.588235,0.349020,0.258824)
rgb=(0.580392,0.349020,0.298039)
rgb=(0.635294,0.349020,0.176471)
rgb=(0.552941,0.360784,0.309804)
rgb=(0.560784,0.364706,0.282353)
rgb=(0.568627,0.368627,0.262745)
rgb=(0.537255,0.372549,0.372549)
rgb=(0.643137,0.380392,0.015686)
rgb=(0.541176,0.392157,0.231373)
rgb=(0.552941,0.376471,0.352941)
rgb=(0.501961,0.388235,0.392157)
rgb=(0.435294,0.400000,0.478431)
rgb=(0.439216,0.403922,0.443137)
rgb=(0.384314,0.415686,0.501961)
rgb=(0.411765,0.411765,0.470588)
rgb=(0.470588,0.403922,0.403922)
rgb=(0.611765,0.400000,0.043137)
rgb=(0.356863,0.431373,0.556863)
rgb=(0.580392,0.407843,0.164706)
rgb=(0.509804,0.407843,0.454902)
rgb=(0.580392,0.400000,0.345098)
rgb=(0.572549,0.431373,0.082353)
rgb=(0.662745,0.372549,0.423529)
rgb=(0.682353,0.388235,0.282353)
rgb=(0.572549,0.419608,0.349020)
rgb=(0.592157,0.415686,0.388235)
rgb=(0.533333,0.458824,0.200000)
rgb=(0.541176,0.470588,0.066667)
rgb=(0.596078,0.458824,0.023529)
rgb=(0.592157,0.423529,0.423529)
rgb=(0.435294,0.462745,0.549020)
rgb=(0.607843,0.431373,0.372549)
rgb=(0.556863,0.439216,0.447059)
rgb=(0.670588,0.450980,0.000000)
rgb=(0.462745,0.462745,0.521569)
rgb=(0.490196,0.458824,0.494118)
rgb=(0.521569,0.454902,0.450980)
rgb=(0.588235,0.443137,0.403922)
rgb=(0.600000,0.447059,0.341176)
rgb=(0.584314,0.447059,0.431373)
rgb=(0.603922,0.447059,0.384314)
rgb=(0.596078,0.466667,0.294118)
rgb=(0.588235,0.501961,0.000000)
rgb=(0.635294,0.490196,0.007843)
rgb=(0.533333,0.498039,0.317647)
rgb=(0.603922,0.474510,0.400000)
rgb=(0.607843,0.474510,0.454902)
rgb=(0.615686,0.478431,0.435294)
rgb=(0.647059,0.466667,0.490196)
rgb=(0.611765,0.490196,0.498039)
rgb=(0.525490,0.513725,0.549020)
rgb=(0.564706,0.505882,0.545098)
rgb=(0.490196,0.521569,0.611765)
rgb=(0.576471,0.509804,0.509804)
rgb=(0.517647,0.521569,0.584314)
rgb=(0.427451,0.541176,0.674510)
rgb=(0.600000,0.505882,0.549020)
rgb=(0.541176,0.600000,0.023529)
rgb=(0.694118,0.529412,0.458824)
rgb=(0.643137,0.549020,0.549020)
rgb=(0.556863,0.568627,0.627451)
rgb=(0.639216,0.549020,0.596078)
rgb=(0.600000,0.556863,0.635294)
rgb=(0.603922,0.564706,0.552941)
rgb=(0.647059,0.564706,0.443137)
rgb=(0.541176,0.576471,0.666667)
rgb=(0.580392,0.572549,0.600000)
rgb=(0.623529,0.564706,0.603922)
rgb=(0.666667,0.600000,0.623529)
rgb=(0.658824,0.600000,0.654902)
rgb=(0.643137,0.611765,0.623529)
rgb=(0.705882,0.596078,0.603922)
rgb=(0.600000,0.623529,0.686275)
rgb=(0.705882,0.603922,0.572549)
rgb=(0.580392,0.627451,0.733333)
rgb=(0.615686,0.627451,0.662745)
rgb=(0.635294,0.623529,0.694118)
rgb=(0.698039,0.615686,0.603922)
rgb=(0.662745,0.619608,0.698039)
rgb=(0.698039,0.658824,0.670588)
rgb=(0.694118,0.658824,0.749020)
rgb=(0.670588,0.670588,0.709804)
rgb=(0.670588,0.674510,0.756863)
rgb=(0.654902,0.682353,0.792157)
rgb=(0.745098,0.701961,0.717647)
rgb=(0.717647,0.721569,0.772549)
rgb=(0.709804,0.721569,0.835294)
rgb=(0.694118,0.745098,0.850980)
rgb=(0.717647,0.776471,0.898039)
}}

\begin{axis}[view={0}{90},axis equal,enlargelimits=false]
    \addplot3 [
        matrix plot,
        colormap access=direct,
        mesh/rows=128,
        mesh/ordering=colwise,
    ] table {legoplane.dat};
\end{axis}
\end{tikzpicture}
\end{codeexample}
    %
    文件 |legoplane_cm.tex|、|legopolane.dat| 和 |legoplane.tiff| 随
    \PGFPlots{} 发布，可从 CTAN 下载。将表格数据解码为 |mesh| 的键
    特定于我们脚本的 Matlab 导出格式。

    注意，从坐标渲染图像在效率上远不如 |\includegraphics|，
    就
    \begin{itemize}
        \item computational time required by \PGFPlots{},
        \item 生成的 |pdf| 文档中的空间消耗，
        \item 以及在查看器中渲染生成的 |pdf| 的时间。
    \end{itemize}
    %
    根据用例，\verbpdfref{\addplot graphics} 可能有助于减少排版时间。


    \paragraph{更多细节：}

    \begin{itemize}
        \item 如果颜色映射中有 $m$ 种颜色，且颜色数据落在 $[0,m-1]$ 之外，
            它将被裁剪到第一种或最后一种颜色。
        \item 如果颜色数据是实数，它将被截断到下一个较小整数。
        \item 如果 |colormap access=direct| 与 |shader|$\neq$|flat corner|
            组合，则进行颜色插值。但是，颜色插值不是按标量
            |point meta| 进行的（如 |colormap access=map| 那样）。
            而是在其自身颜色空间中的\emph{结果}颜色上进行。
        \item 注意，在此上下文中不支持 |patch refines| 或
            |patch to triangles|（在撰写本文时，它将插值颜色索引，
            导致定义不良的行为）。
    \end{itemize}

    另见相关键 |color of colormap| 和 |index of colormap|。

    选项 |colormap access=|\declaretext{piecewise linear} 是
    |colormap access=map| 的别名。它与 |color of colormap| 相关。
    \index{color of colormap}


    \subsubsection{使用分段常数插值映射元数据}

    \paragraph{注意：查看器限制}

    以下段落使用 \textsc{pdf} 功能，某些免费查看器
    （包括 |pdf.js| 和 |evince|）不支持。如果看到意外结果，
    请改用合适的查看器。

    选项 |colormap access=|\declaretext{piecewise constant} 与
    |colormap access=map| 非常相似：它接受相同的 |point meta| 值作为输入，
    线性映射到 $[0,1000]$，如果 |point meta| 值为 |point meta min| 则使用
    |colormap| 的第一个颜色，如果 |point meta| 值为 |point meta max|
    则使用 |colormap| 的最后一个颜色。但是，落在两个 |colormap| 条目
    \emph{之间} 的 |point meta| 值使用\emph{分段常数}插值，即就像
    |const plot|：插值颜色类似于区间\emph{第一个}颜色的值。因此，
    结果包含的颜色最多与 |colormap| 提供的一样多：
    %
\begin{codeexample}[]
% requires \usepgfplotslibrary{colorbrewer}
\begin{tikzpicture}
\begin{axis}[
    colormap access=piecewise constant,
    colormap/RdPu-6,
    colorbar horizontal,
]
    \addplot3 [
        surf,
    ] {x*y};
\end{axis}
\end{tikzpicture}
\end{codeexample}
    %
    选项 |colormap access=piecewise constant| 与所有可用 |shader| 兼容，
    但使用 |shader=interp| 可获得最佳结果，因为这会避免网格边界处的
    矩形伪影。
    %
\begin{codeexample}[]
% requires \usepgfplotslibrary{colorbrewer}
\begin{tikzpicture}
\begin{axis}[
    colormap access=piecewise constant,
    colormap/RdPu-6,
    colorbar horizontal,
]
    \addplot3 [
        surf,
        shader=interp,
    ] {x*y};
\end{axis}
\end{tikzpicture}
\end{codeexample}

    选项 |colormap access=piecewise constant| 需要具有合适颜色数量的
    |colormap|。这与 |colormap access=map|（或 |piecewise linear|）非常不同，
    后者中 |colormap| 的实际颜色数量在结果图中不可见。为获得所需效果，
    可能需要具有不同样本数量的相同 |colormap|。为此，可以使用基于
    |samples of colormap| 及其变体的特殊 |colormap| 语法：

\begin{codeexample}[]
\begin{tikzpicture}
\begin{axis}[
    colormap access=piecewise constant,
    colormap={my sampled version of viridis}{
        samples of colormap=(10 of viridis)
    },
    colorbar horizontal,
]
    \addplot3 [
        surf,
        shader=interp,
    ] {x*y};
\end{axis}
\end{tikzpicture}
\end{codeexample}
    %
    该示例使用现有 |colormap name=viridis| 并选择其 $10$ 个样本。
    结果称为 |my sampled version of viridis|，用作``当前颜色映射''。
    有关此过程的详情，请阅读第
    ~\ref{sec:pgfplots:colormaps:based:on:others} 节
    （第~\pageref{sec:pgfplots:colormaps:based:on:others} 页）。


    \paragraph{另见}

    绘图处理程序 |contour filled|，它基于此类可视化方案。

    在此上下文中确定颜色的策略也可作为
    |const color of colormap| 使用。\index{const color of colormap}

    选项 \declareandlabel{piecewise const} 和 \declaretext{const} 是
    |piecewise constant| 的别名。
\end{pgfplotskey}
